% TOOL:    airfoil-drag-aircraft-performance.tex - IEEE-style two-column A4 layout (generated by build.py)
% FAMILY:  technical-papers / airfoil-drag-aircraft-performance
% VERSION: 1.0.0
% DATE:    2026-11-28
% CHAT:    Airfoil drag characteristics and aircraft/rotorcraft performance (Aerodynamics Related project)
% CHANGES: 1.0.0 - first version: Times text and maths, Roman section numbers, IEEE table and figure captions
\documentclass[10pt,a4paper,twocolumn]{article}
\usepackage[a4paper,top=19mm,bottom=24mm,left=15mm,right=15mm,columnsep=5mm,footskip=10mm,headsep=5mm,headheight=12pt]{geometry}
\usepackage[T1]{fontenc}
\usepackage[utf8]{inputenc}
\usepackage{textcomp}
\usepackage{amsmath,amssymb}
\usepackage{mathptmx}
\usepackage{newunicodechar}
\newunicodechar{−}{\ensuremath{-}}
\newunicodechar{≈}{\ensuremath{\approx}}
\newunicodechar{→}{\ensuremath{\rightarrow}}
\newunicodechar{Δ}{\ensuremath{\Delta}}
\newunicodechar{⁻}{\ensuremath{^{-}}}
\newunicodechar{⁷}{\ensuremath{^{7}}}
\newunicodechar{⁶}{\ensuremath{^{6}}}
\newunicodechar{₂}{\ensuremath{_{2}}}
\usepackage{graphicx,booktabs,tabularx,array,calc,xcolor}
\usepackage[font=footnotesize,labelformat=empty,justification=justified,singlelinecheck=true,skip=4pt]{caption}
\usepackage{titlesec}
\usepackage{fancyhdr}
\usepackage{stfloats}
\usepackage[hidelinks,unicode]{hyperref}
\providecommand{\tightlist}{\setlength{\itemsep}{0pt}\setlength{\parskip}{0pt}}
\setlength{\emergencystretch}{2em}
\raggedbottom
\renewcommand\thesection{\Roman{section}}
\renewcommand\thesubsection{\Alph{subsection}}
\titleformat{\section}{\centering\normalfont\scshape}{\thesection.}{0.5em}{}
\titleformat{\subsection}{\normalfont\itshape}{\thesubsection.}{0.5em}{}
\titlespacing*{\section}{0pt}{1.6ex plus .4ex minus .2ex}{0.8ex plus .2ex}
\titlespacing*{\subsection}{0pt}{1.0ex plus .3ex minus .2ex}{0.5ex}
\setlength{\parindent}{1em}
\setlength{\parskip}{0pt}
\setlength{\abovedisplayskip}{5pt}\setlength{\belowdisplayskip}{5pt}
\setlength{\textfloatsep}{10pt plus 2pt minus 2pt}
\setlength{\dbltextfloatsep}{10pt plus 2pt minus 2pt}
\renewcommand{\topfraction}{0.9}\renewcommand{\dbltopfraction}{0.9}
\renewcommand{\textfraction}{0.08}\renewcommand{\floatpagefraction}{0.8}\renewcommand{\dblfloatpagefraction}{0.8}
\pagestyle{fancy}\fancyhf{}
\fancyhead[C]{\footnotesize\itshape From Section Drag to Mission Performance}
\fancyfoot[C]{\footnotesize\thepage}
\renewcommand{\headrulewidth}{0pt}
\fancypagestyle{plain}{\fancyhf{}\fancyfoot[C]{\footnotesize\thepage}\renewcommand{\headrulewidth}{0pt}}
\hypersetup{pdftitle={From Section Drag to Mission Performance: Airfoil Choice Across Fixed-Wing and Rotary Aircraft},pdfauthor={K Vijay Anand}}
\begin{document}
\twocolumn[{%
\begin{center}
{\fontsize{22}{26}\selectfont From Section Drag to Mission Performance:\\[2pt] Airfoil Choice Across Fixed-Wing and Rotary Aircraft\par}
\vspace{10pt}
{\large K Vijay Anand\par}
\vspace{2pt}
{\small November 2026\par}
\vspace{14pt}
\end{center}
}]
\thispagestyle{plain}
\begingroup\small\bfseries

\textbf{Abstract}---The drag characteristics of the two-dimensional
airfoil section---its minimum profile drag, the width and position of
its low-drag range, its behaviour with Reynolds number, transition, Mach
number and surface condition---propagate into every first-order measure
of aircraft performance: range, endurance, climb, ceiling, manoeuvre,
hover capability and payload. This paper traces that propagation for the
main classes of fixed-wing aircraft and rotorcraft: subsonic transports,
turboprops, sailplanes and long-endurance unmanned aircraft, small
low-Reynolds-number drones, combat aircraft, helicopters, tiltrotors,
compounds and multicopters. A common framework is built from a simple
but calibrated section-drag model (mixed laminar--turbulent skin
friction with a thickness form factor), the parabolic drag polar, the
Breguet range and endurance relations, the Korn drag-divergence
relation, linear supersonic theory and blade-element momentum theory for
rotors and propellers. Ten worked examples (A--J), with supporting
turboprop and propeller cases, quantify the sensitivities. For a generic
narrow-body transport, a 10 \% reduction in wing profile drag reduces
cruise fuel by 2.3 \% on a 2 500 km sector, and a hybrid-laminar wing
that cuts section drag by about 29 \% saves 5.3 \%. For a utility
helicopter in hot-and-high hover, a 10 \% lower blade drag polar adds
about 47 kg of hover-capable mass, while a 25 \% degradation from
leading-edge erosion removes about 127 kg. For a small quadcopter,
raising the rotor profile drag coefficient from 0.015 to 0.035 lowers
figure of merit from 0.69 to 0.59 and endurance by 15 \%. The results
show that the \emph{same} section drag change has very different
consequences depending on whether the vehicle is range-, endurance-,
power- or Mach-limited, and that for rotorcraft and supersonic aircraft
it is often the drag \emph{rise} (with lift coefficient or Mach number),
rather than the minimum drag, that sets the performance boundary. The
paper closes with a synthesis table, design guidance and a note on the
prediction tools---XFOIL-type viscous--inviscid codes, RANS and
wind-tunnel data---needed to quantify these effects.

\par\endgroup\smallskip

\begingroup\small\bfseries

\textbf{\emph{Index Terms}}---airfoil, profile drag, drag polar, laminar
flow, transonic drag divergence, wave drag, Breguet range, endurance,
helicopter performance, figure of merit, propeller efficiency, low
Reynolds number, UAV.

\par\endgroup\smallskip

\hypertarget{nomenclature}{%
\section*{Nomenclature}\label{nomenclature}}
\addcontentsline{toc}{section}{Nomenclature}

{\footnotesize\setlength{\tabcolsep}{4pt}\noindent \begin{tabularx}{\columnwidth}{l>{\raggedright\arraybackslash}X}
\toprule
\textbf{Symbol} & \textbf{Meaning} \\
\midrule
\(A\) & Aspect ratio (wing); disk area \(\pi R^2\) (rotor) \\
\(a\) & Speed of sound \\
\(c\) & Chord; thrust-specific fuel consumption (jet) \\
\(c_p\) & Power-specific fuel consumption \\
\(c_d\), \(c_l\) & Section drag and lift coefficients \\
\(C_D\), \(C_L\), \(C_{D0}\) & Aircraft drag, lift and zero-lift drag coefficients \\
\(C_f\) & Mean skin-friction coefficient \\
\(C_T\) & Rotor thrust coefficient \(T/(\rho A(\Omega R)^2)\) \\
\(e\) & Oswald span-efficiency factor \\
\(f\) & Equivalent flat-plate area \\
\(FM\) & Rotor figure of merit \\
\(K\) & Induced-drag factor \(1/(\pi A e)\) \\
\(M\), \(M_{cr}\), \(M_{DD}\) & Mach number; critical and drag-divergence Mach numbers \\
\(q\) & Dynamic pressure \(\tfrac12\rho V^2\) \\
\(R\) & Range; rotor radius \\
\(Re\) & Reynolds number based on chord \\
\(S\) & Wing reference area \\
\(t/c\), \(\tau\) & Thickness-to-chord ratio \\
\(T\), \(W\), \(V\) & Thrust, weight, flight speed \\
\(x_{tr}\) & Transition location \\
\(\beta\) & \(\sqrt{M^2-1}\) \\
\(\eta_p\) & Propulsive efficiency \\
\(\kappa\), \(\kappa_A\) & Induced-power factor (rotor); Korn technology factor \\
\(\Lambda\) & Sweep angle \\
\(\mu\) & Rotor advance ratio \(V/\Omega R\) \\
\(\rho\) & Air density \\
\(\sigma\) & Rotor solidity \\
\(\phi\) & Blade-element inflow angle \\
\(\Omega R\) & Rotor tip speed \\
\bottomrule
\end{tabularx}\par}\medskip

\hypertarget{introduction}{%
\section{Introduction}\label{introduction}}

Every aircraft converts stored energy into distance, time aloft or
useful work against gravity, and drag is the principal tax on that
conversion. Within the total drag of an aircraft, the contribution of
the lifting surfaces' \emph{sections}---the airfoils---occupies a
special place. It is the part of the drag that the designer selects
early, often from a catalogue or a code such as XFOIL {[}1{]}, and it
then constrains wing or blade thickness, structural weight, maximum
lift, buffet margin, propulsive efficiency and operating Mach number for
the life of the design.

The influence of airfoil drag is far from uniform across aircraft types.
On a jet transport the wing section is chosen chiefly for its
\emph{drag-divergence Mach number} at a given thickness and lift
coefficient, and its subsonic profile drag accounts for perhaps a third
of zero-lift drag {[}2{]}, {[}3{]}. On a sailplane the profile drag of a
laminar airfoil is the dominant single term in the glide polar, and
losing laminar flow to insect debris can cost more glide performance
than any other operational factor {[}4{]}. On a helicopter, profile
power is only 10--30 \% of the total in hover but it is the term that
grows fastest at high forward speed, and the airfoil's drag rise with
Mach number on the advancing blade and with lift coefficient on the
retreating blade defines the flight envelope {[}5{]}, {[}6{]}, {[}7{]}.
On small drones the airfoil operates at Reynolds numbers where laminar
separation bubbles can double the section drag {[}8{]}, {[}9{]},
{[}10{]}.

This paper asks a simple question: \emph{how does a change in airfoil
drag characteristics change the performance of the whole vehicle, and
why does the answer differ between aircraft types?} It does so by:

\begin{enumerate}
\def\labelenumi{\arabic{enumi}.}
\tightlist
\item
  reviewing the physics of section drag and the parameters that control
  it (Section II);
\item
  setting out the performance relations that connect section drag to
  vehicle metrics, and deriving their sensitivities (Section III);
\item
  examining, with worked numerical examples, subsonic jet transports
  (Section IV), turboprops, light aircraft and propellers (Section V),
  sailplanes and long-endurance aircraft (Section VI), small
  low-Reynolds-number aircraft (Section VII), combat aircraft (Section
  VIII), helicopters (Section IX), and tiltrotors, compounds and
  multicopters (Section X);
\item
  considering second-order couplings through pitching moment and
  propulsion integration (Section XI) and illustrative case studies
  (Section XII);
\item
  discussing in-service degradation (Section XIII), prediction methods
  (Section XIV) and multi-point airfoil selection (Section XV); and
\item
  drawing the results together in a synthesis and design guidance
  (Section XVI) and conclusions (Section XVII), with the example
  parameters and analytical sensitivities in two appendices.
\end{enumerate}

All numerical examples use generic, representative parameters drawn from
open textbook data; they are not models of any specific aircraft. They
are intended to show \emph{orders of magnitude and trends}, and are
computed with a single open script so that each number can be reproduced
and varied.

\hypertarget{historical-perspective}{%
\subsection{Historical perspective}\label{historical-perspective}}

The link between airfoil drag and aircraft performance has driven
airfoil development since the earliest days of flight. The thin, highly
cambered sections of the first aeroplanes were chosen for lift, not
drag; their poor \(c_l/c_d\) and abrupt stall contributed to the low
speeds and short ranges of the period. The Göttingen and NACA
wind-tunnel programmes of the 1920s and 1930s showed that thicker,
rounded-nose sections had \emph{lower} drag as well as higher maximum
lift, and allowed cantilever monoplane wings whose lower parasite drag
transformed cruise speed and range. The NACA four- and five-digit
families systematised this knowledge {[}11{]}.

The next step was deliberate control of transition. The NACA 6-series
sections of the late 1930s and 1940s were shaped to hold a favourable
pressure gradient over a large part of the chord, producing the low-drag
bucket described in Section II. In practice the laminar extent achieved
on production aircraft of that era was limited by surface quality, and
the main lasting benefit was often the sections' higher critical Mach
number---an early example of a drag characteristic other than minimum
drag proving decisive {[}11{]}, {[}12{]}.

The jet age moved the emphasis to compressibility. Sweep raised
drag-divergence Mach numbers but at a cost in structural weight and
low-speed handling. Whitcomb's supercritical airfoil of the mid-1960s
{[}13{]}, {[}14{]} showed that careful shaping of the upper-surface
supersonic region could raise \(M_{DD}\) at a given thickness, or allow
a thicker section at a given \(M_{DD}\); it is now the basis of every
transport wing. In parallel, sailplane designers such as Wortmann and
Eppler, using early computer methods, pushed natural laminar flow to its
practical limits at moderate Reynolds number {[}15{]}, {[}4{]}, and
NASA's NLF programme of the 1980s brought such sections to general
aviation {[}16{]}.

Rotorcraft followed a somewhat separate path. For decades helicopter
blades used symmetric NACA 0012-type sections because of their small
pitching moment. From the 1970s, dedicated rotorcraft airfoil families
were developed to combine high \(c_{l,\max}\) at low Mach number (for
the retreating blade) with high drag-divergence Mach number at low lift
(for the advancing blade) while keeping pitching moments small {[}17{]},
{[}5{]}. Over the same period, the rise of small unmanned aircraft
revived interest in the low-Reynolds-number regime that had been studied
for model aircraft, and dedicated low-Reynolds-number airfoil data sets
were produced {[}9{]}, {[}10{]}.

Loftin's survey of aircraft performance evolution {[}12{]} shows the
outcome of this history: maximum lift-to-drag ratios rose from below 10
for the biplanes of the First World War to around 15--20 for modern
transports and above 40 for modern sailplanes. Airfoil drag reduction
was one strand in that progress, together with higher aspect ratio,
cleaner configurations and better propulsion. The present paper looks at
the same history from the opposite direction: given a change in the
airfoil, how much does the vehicle gain?

\hypertarget{airfoil-drag-components-and-controlling-parameters}{%
\section{Airfoil Drag: Components and Controlling
Parameters}\label{airfoil-drag-components-and-controlling-parameters}}

\hypertarget{decomposition-of-section-drag}{%
\subsection{Decomposition of section
drag}\label{decomposition-of-section-drag}}

The drag of a two-dimensional airfoil section per unit span is expressed
as a coefficient \(c_d = D'/(q_\infty c)\), where \(q_\infty\) is the
dynamic pressure and \(c\) the chord. It is conventionally divided into
three parts:

\[c_d = c_{d,f} + c_{d,p} + c_{d,w} \tag{1}\]

The \emph{skin-friction drag} \(c_{d,f}\) is the integrated wall shear
stress. The \emph{pressure} or \emph{form drag} \(c_{d,p}\) arises
because the boundary layer displaces the outer flow and, near the
trailing edge, prevents full pressure recovery; it rises sharply once
the flow separates. Together these make up the \emph{profile drag}. The
\emph{wave drag} \(c_{d,w}\) appears only when the local flow becomes
supersonic and shock waves form---above the critical Mach number in
transonic flow, and everywhere in supersonic flight {[}18{]}, {[}19{]}.

In the language of performance, a ``drag count'' is
\(\Delta c_d = 0.0001\). A modern turbulent transport wing section at
cruise has a profile drag of roughly 55--70 counts; a good laminar
sailplane section 45--60 counts; a helicopter blade section 80--130
counts averaged over the operating range; and a small-drone airfoil at a
Reynolds number of \(10^5\) commonly 120--250 counts {[}5{]}, {[}11{]},
{[}9{]}.

\hypertarget{reynolds-number-and-transition}{%
\subsection{Reynolds number and
transition}\label{reynolds-number-and-transition}}

Skin friction depends first on Reynolds number \(Re = \rho V c/\mu\) and
on whether the boundary layer is laminar or turbulent. A laminar layer
has a skin-friction coefficient several times lower than a turbulent one
at the same Reynolds number. A widely used engineering estimate of the
mean skin friction of a surface with transition at \(x_{tr}/c\) combines
the Blasius laminar and Prandtl--Schlichting turbulent results {[}19{]}:

\[\bar C_f = C_{f,t}(Re) - \frac{x_{tr}}{c}\left[C_{f,t}(Re_{tr}) - C_{f,l}(Re_{tr})\right] \tag{2}\]

with \(C_{f,l}=1.328/\sqrt{Re}\), \(C_{f,t}=0.455/(\log_{10}Re)^{2.58}\)
and \(Re_{tr}=Re\,x_{tr}/c\). Pressure drag of an attached-flow section
is accounted for by a thickness form factor, giving a section profile
drag of {[}20{]}, {[}2{]}

\[c_d \approx \left(\bar C_{f,\text{upper}}+\bar C_{f,\text{lower}}\right)\left[1+2\frac{t}{c}+60\left(\frac{t}{c}\right)^4\right] \tag{3}\]

Fig. 1 shows (3) for a 12 \% thick section. Two effects dominate. First,
at a fixed transition position, drag falls with Reynolds number---by
about 45 \% between \(Re=10^6\) and \(3\times10^7\) for a fully
turbulent section---so small aircraft and rotor blades pay a scale
penalty. Second, moving transition from near the leading edge to 60 \%
chord reduces drag by 40--55 \% at every Reynolds number. This is the
prize of \emph{natural laminar flow} (NLF) airfoils such as the NACA
6-series and the NLF(1)-0416 {[}11{]}, {[}16{]}, and of \emph{laminar
flow control} (LFC) {[}21{]}.

\begin{figure}[!t]
\centering
\includegraphics[width=\columnwidth]{fig/f1_cd_transition.pdf}
\caption{Fig. 1. Section profile drag of a 12 \% thick airfoil from (2)--(3) as a
function of transition location, for Reynolds numbers typical of drones
(\(10^6\)), light aircraft and rotor blades (\(3\times10^6\)), and
transports (\(10^7\)--\(3\times10^7\)).}
\end{figure}

Table I lists values from the model. The model is deliberately simple:
it ignores the pressure-gradient effects that real airfoils exploit, but
it reproduces the drag of well-designed attached-flow sections to within
10--15 \% and, more importantly, captures the \emph{relative} effect of
Reynolds number and transition that the later examples rely on.

\begin{table}[!t]
\centering
{\footnotesize TABLE I\par\textsc{Model section profile drag \(c_d\) (counts) of a 12 \% thick
section, equations (2)--(3).}\par}\vspace{4pt}
{\footnotesize\setlength{\tabcolsep}{4pt}\renewcommand{\arraystretch}{1.12}
\begin{tabular}{lrrr}
\toprule
\textbf{\(Re\)} & \textbf{Turbulent (\(x_{tr}\)=0.05)} & \textbf{\(x_{tr}\)=0.40} & \textbf{\(x_{tr}\)=0.60} \\
\midrule
\(1\times10^6\) & 109 & 80 & 64 \\
\(3\times10^6\) & 88 & 61 & 47 \\
\(1\times10^7\) & 71 & 47 & 34 \\
\(3\times10^7\) & 60 & 38 & 27 \\
\bottomrule
\end{tabular}
}
\end{table}

Whether transition actually occurs late depends on the pressure
distribution (a favourable gradient stabilises the layer), on surface
waviness, steps, gaps and roughness, on free-stream turbulence, on
leading-edge sweep (which introduces crossflow and attachment-line
instabilities) and on contamination by insects, rain or ice. In
practice, transition is predicted with the \(e^N\) method, in which the
amplification of Tollmien--Schlichting waves is tracked until a critical
factor---typically \(N=9\) in flight and lower in a turbulent wind
tunnel---is reached {[}1{]}, {[}19{]}.

\hypertarget{variation-with-lift-coefficient-the-drag-bucket}{%
\subsection{Variation with lift coefficient: the drag
bucket}\label{variation-with-lift-coefficient-the-drag-bucket}}

The section drag polar \(c_d(c_l)\) is not flat. Laminar sections show a
\emph{drag bucket}: a range of \(c_l\) over which favourable gradients
on both surfaces hold transition aft and drag is low. Outside the
bucket, transition jumps forward on one surface and drag rises abruptly.
Conventional turbulent sections have a broader, shallower minimum. At
high \(c_l\) all sections show a steep rise as trailing-edge separation
develops ahead of stall. For performance, the location of the bucket
relative to the vehicle's operating \(c_l\)---cruise, climb, loiter or,
for a rotor, the azimuthally varying blade \(c_l\)---matters as much as
the minimum drag itself.

Fig. 2 illustrates the three characteristic shapes. The laminar section
has the lowest drag inside its bucket but the sharpest rise outside it;
the turbulent section has higher minimum drag but a broader, more
forgiving polar; and the low-Reynolds-number section has both higher
minimum drag and a narrower useful range, because laminar separation
bubbles add drag at the ends of its operating range. Which of these is
``best'' depends entirely on where the vehicle operates on the \(c_l\)
axis and how widely that operating point moves.

\begin{figure}[!t]
\centering
\includegraphics[width=\columnwidth]{fig/f0b_polars.pdf}
\caption{Fig. 2. Illustrative section drag polars (stylised from published
trends) for a natural-laminar-flow section, a conventional turbulent
section and a low-Reynolds-number section.}
\end{figure}

\hypertarget{compressibility-and-wave-drag}{%
\subsection{Compressibility and wave
drag}\label{compressibility-and-wave-drag}}

As free-stream Mach number \(M_\infty\) increases, the suction peak on
the upper surface reaches sonic speed at the critical Mach number
\(M_{cr}\). A little above it, a shock forms, wave drag appears and the
shock can trigger boundary-layer separation. The \emph{drag-divergence
Mach number} \(M_{DD}\), usually defined where \(dc_d/dM = 0.1\), marks
the practical limit for efficient cruise. \(M_{DD}\) falls with
thickness and lift coefficient and rises with sweep and with
\emph{supercritical} section design, which shapes the upper surface to
decelerate the supersonic pocket gently, with a weak or no shock
{[}13{]}, {[}14{]}. In supersonic flight wave drag is present at all
lift coefficients and scales with the square of thickness ratio (Section
VIII).

\hypertarget{profile-drag-from-the-wake-the-squireyoung-relation}{%
\subsection{Profile drag from the wake: the Squire--Young
relation}\label{profile-drag-from-the-wake-the-squireyoung-relation}}

Profile drag can be obtained from the state of the boundary layer at the
trailing edge. Far downstream, the drag per unit span equals the
momentum deficit of the wake, \(c_d = 2\theta_\infty/c\), where
\(\theta_\infty\) is the far-wake momentum thickness. Squire and Young
related this to the trailing-edge values of momentum thickness
\(\theta_{TE}\), shape factor \(H_{TE}\) and edge velocity \(U_{e,TE}\)
{[}19{]}:

\[c_d = \frac{2\theta_{TE}}{c}\left(\frac{U_{e,TE}}{U_\infty}\right)^{(H_{TE}+5)/2} \tag{4}\]

The relation, or its equivalent, is used by viscous--inviscid codes such
as XFOIL {[}1{]} to evaluate drag, and it makes the physical levers
explicit. Drag rises with the momentum thickness accumulated along each
surface---which is small for a laminar layer and large for a turbulent
one---and with the shape factor at the trailing edge, which grows as the
layer approaches separation. A section that keeps the layer laminar for
longer, and brings it to the trailing edge with a moderate adverse
gradient so that \(H_{TE}\) stays low, has low drag. A section with a
steep pressure recovery, or one operating near stall, pays through both
\(\theta_{TE}\) and \(H_{TE}\).

\hypertarget{transition-on-swept-wings-and-rotating-blades}{%
\subsection{Transition on swept wings and rotating
blades}\label{transition-on-swept-wings-and-rotating-blades}}

On an unswept wing in a quiet free stream, transition is governed mainly
by Tollmien--Schlichting waves, which a favourable pressure gradient
suppresses. Sweep introduces two further mechanisms {[}21{]}.
\emph{Crossflow instability} arises because the curved streamlines in a
swept boundary layer create a velocity component normal to the inviscid
streamline; it is \emph{amplified} by the strong favourable gradient
near the leading edge that stabilises Tollmien--Schlichting waves, so
the two must be balanced. \emph{Attachment-line contamination} occurs
when turbulence from the fuselage junction travels along the
leading-edge attachment line; above a critical attachment-line Reynolds
number it cannot be stopped without devices such as a leading-edge bump
or suction. Together these mechanisms make laminar flow increasingly
difficult as sweep and Reynolds number increase, which is why natural
laminar flow is common on sailplanes and business jets but has not yet
entered service on large transport wings.

Rotor blades add rotation and unsteadiness. Each blade section passes
through the wake of the preceding blade and sees periodic changes in
incidence and Mach number, so transition moves back and forth around the
azimuth. In practice, outboard rotor sections are often treated as
largely turbulent for performance prediction, although measurable
laminar runs exist on smooth composite blades in hover.

\hypertarget{from-section-drag-to-wing-and-blade-drag}{%
\subsection{From section drag to wing and blade
drag}\label{from-section-drag-to-wing-and-blade-drag}}

A finite wing adds induced drag, the cost of producing lift with a
finite span, and interference and junction drag. The section
characteristics enter the three-dimensional drag in two ways: through
the \emph{zero-lift} or \emph{parasite} contribution of the wing,
approximately \(c_d\) times the exposed-to-reference area ratio, and
through the \emph{lift-dependent} part of profile drag, which the
parabolic polar absorbs into the Oswald factor \(e\). For a rotor, the
section polar is integrated over radius and azimuth, with each element
operating at its own Reynolds number, Mach number and \(c_l\) (Section
IX). Table II summarises which drag mechanism dominates for each vehicle
class and therefore which airfoil characteristic matters most. Fig. 3
places the same classes on a Reynolds--Mach map; the spread covers four
decades of Reynolds number and two of Mach number, which is why no
single airfoil family, and no single drag figure, serves all of them.

\begin{figure}[!t]
\centering
\includegraphics[width=\columnwidth]{fig/f0_re_mach_map.pdf}
\caption{Fig. 3. Approximate operating ranges of chord Reynolds number and Mach
number for the vehicle classes considered.}
\end{figure}

\begin{table*}[!t]
\centering
{\footnotesize TABLE II\par\textsc{Where airfoil drag matters, by vehicle class.}\par}\vspace{4pt}
{\footnotesize\setlength{\tabcolsep}{4pt}\renewcommand{\arraystretch}{1.12}
\begin{tabularx}{\textwidth}{l>{\raggedright\arraybackslash}Xl>{\raggedright\arraybackslash}X>{\raggedright\arraybackslash}X}
\toprule
\textbf{Vehicle class} & \textbf{Typical \(Re\)} & \textbf{Typical \(M\)} & \textbf{Dominant section characteristic} & \textbf{Main performance lever} \\
\midrule
Jet transport & \(1\)--\(4\times10^7\) & 0.75--0.85 & \(M_{DD}\) at given \(t/c\), \(c_l\); turbulent \(c_{d,\min}\) & Range, fuel burn, wing weight \\
Turboprop / GA & \(3\times10^6\)--\(10^7\) & 0.2--0.6 & \(c_{d,\min}\), bucket, NLF extent & Range, climb, prop efficiency \\
Sailplane / HALE & \(5\times10^5\)--\(3\times10^6\) & \textless0.2 & Laminar bucket, low-\(c_l\) and high-\(c_l\) drag & Glide ratio, sink, endurance \\
Small UAV & \(5\times10^4\)--\(5\times10^5\) & \textless0.1 & Laminar separation bubble, \(c_l^{3/2}/c_d\) & Endurance (battery) \\
Combat aircraft & \(2\times10^7\)--\(10^8\) & 0.6--2.0 & Thin-section wave drag, LE suction at high \(c_l\) & Sustained turn, supersonic acceleration \\
Helicopter & \(2\)--\(8\times10^6\) & 0.3--0.85 & \(c_{d0}\), drag rise with \(M\) (advancing) and \(c_l\) (retreating) & Hover payload, speed, range \\
Multicopter & \(3\times10^4\)--\(2\times10^5\) & \textless0.35 & Low-\(Re\) \(c_{d0}\), \(c_{l,\max}\) & Hover endurance \\
\bottomrule
\end{tabularx}
}
\end{table*}

\hypertarget{linking-section-drag-to-vehicle-performance}{%
\section{Linking Section Drag to Vehicle
Performance}\label{linking-section-drag-to-vehicle-performance}}

\hypertarget{the-drag-polar-and-maximum-lift-to-drag-ratio}{%
\subsection{The drag polar and maximum lift-to-drag
ratio}\label{the-drag-polar-and-maximum-lift-to-drag-ratio}}

For performance analysis the aircraft drag is usually approximated by
the parabolic polar

\[C_D = C_{D0} + K C_L^2, \qquad K = \frac{1}{\pi A e} \tag{5}\]

where \(A\) is aspect ratio and \(e\) the span efficiency (Oswald)
factor. \(C_{D0}\) is built up from the parasite drag of each component,
\(C_{D0}\approx\sum C_{f,i}\,FF_i\,Q_i\,S_{wet,i}/S_{ref}\) {[}2{]},
{[}22{]}; the wing term is the one controlled by the airfoil. The
maximum lift-to-drag ratio and the lift coefficient at which it occurs
are

\[\left(\frac{L}{D}\right)_{\max} = \frac{1}{2\sqrt{C_{D0}K}}, \qquad C_L^* = \sqrt{\frac{C_{D0}}{K}} \tag{6}\]

\hypertarget{range}{%
\subsection{Range}\label{range}}

For a jet with thrust-specific fuel consumption \(c\) (per unit time,
weight basis), the Breguet range at constant \(V\), \(L/D\) and \(c\) is

\[R = \frac{V}{c}\,\frac{L}{D}\,\ln\frac{W_i}{W_f} \tag{7}\]

and for a propeller aircraft with power-specific fuel consumption
\(c_p\) and propulsive efficiency \(\eta_p\),

\[R = \frac{\eta_p}{c_p}\,\frac{L}{D}\,\ln\frac{W_i}{W_f} \tag{8}\]

For a fixed mission range \(R\) ending at weight \(W_f\), the fuel
burned is

\[W_{fuel} = W_f\left[\exp\!\left(\frac{R\,c}{V\,(L/D)}\right)-1\right] \tag{9}\]

so fuel depends on \(L/D\) slightly more than inversely: a 1 \%
improvement in \(L/D\) saves a little more than 1 \% of trip fuel
because less fuel need be carried to carry the fuel.

\hypertarget{endurance}{%
\subsection{Endurance}\label{endurance}}

Endurance for a jet is maximised at \((L/D)_{\max}\). For propeller and
electric aircraft, where power rather than thrust is consumed, endurance
scales with \(C_L^{3/2}/C_D\), whose maximum from (5) is

\[\left(\frac{C_L^{3/2}}{C_D}\right)_{\max} = \frac{1}{4}\left(\frac{3}{K}\right)^{3/4} C_{D0}^{-1/4} \tag{10}\]

occurring at \(C_L=\sqrt{3C_{D0}/K}\), where induced drag is three times
the parasite drag.

\hypertarget{sensitivity-exponents}{%
\subsection{Sensitivity exponents}\label{sensitivity-exponents}}

Equations (6) and (10) give a compact and important result. Near the
optimum, \((L/D)_{\max}\propto C_{D0}^{-1/2}\) but
\((C_L^{3/2}/C_D)_{\max}\propto C_{D0}^{-1/4}\). A 10 \% reduction in
\(C_{D0}\) therefore raises the maximum \(L/D\) (and hence range for an
aircraft flying at \(C_L^*\)) by 5.4 \%, but raises maximum propeller or
electric endurance by only 2.7 \%. Endurance aircraft fly at high
\(C_L\), where induced drag dominates; the airfoil influences them less
through minimum drag and more through the \emph{shape} of the section
polar at high \(c_l\), and through enabling high aspect ratio.
Conversely, an aircraft forced to cruise \emph{below} \(C_L^*\)---at
high speed or low altitude---is \emph{more} sensitive to \(C_{D0}\),
because parasite drag then dominates.

\hypertarget{climb-ceiling-and-manoeuvre}{%
\subsection{Climb, ceiling and
manoeuvre}\label{climb-ceiling-and-manoeuvre}}

Specific excess power \(P_s = V(T-D)/W\) governs rate of climb, ceiling
and energy manoeuvrability. At high speed \(D\) is mostly parasite drag;
at the ceiling and in a sustained turn, it is mostly lift-dependent
drag. For a sustained turn at load factor \(n\) with thrust \(T\):

\[n_{sus} = \frac{qS}{W}\sqrt{\frac{T/(qS) - C_{D0}}{K}} \tag{11}\]

In this regime the section characteristic that matters is the
lift-dependent part of profile drag---how much \emph{leading-edge
suction} the section can sustain before separation---rather than
\(c_{d,\min}\).

\hypertarget{the-airfoils-share-of-the-drag-budget}{%
\subsection{The airfoil's share of the drag
budget}\label{the-airfoils-share-of-the-drag-budget}}

A useful rule follows from these relations. For small changes, the
fractional change in a performance metric is approximately the product
of two factors: the \emph{share} of total drag (or power) that the
airfoil's profile drag represents at the limiting condition, and an
\emph{exponent} of order one that depends on the metric---about 1 for
fuel at fixed range, \(\tfrac12\) for maximum \(L/D\), \(\tfrac13\) for
top speed and \(\tfrac23 P_0/P_i\)-type ratios for hover thrust
(Appendix B). Table III lists the shares for the generic vehicles used
in this paper. They range from about 13 \% for a fighter in a sustained
turn to about 32 \% for a sailplane at best glide, which already
explains most of the spread in sensitivities found later.

\begin{table}[!t]
\centering
{\footnotesize TABLE III\par\textsc{Share of total drag or rotor power due to wing or blade
section profile drag in the generic examples.}\par}\vspace{4pt}
{\footnotesize\setlength{\tabcolsep}{4pt}\renewcommand{\arraystretch}{1.12}
\begin{tabularx}{\columnwidth}{>{\raggedright\arraybackslash}X>{\raggedright\arraybackslash}X}
\toprule
\textbf{Vehicle and condition} & \textbf{Airfoil profile share} \\
\midrule
Jet transport, mid-cruise (A) & 22 \% of drag \\
Regional turboprop, cruise & 26 \% of drag \\
Sailplane at \((L/D)_{\max}\) (D) & 32 \% of drag \\
Mini-UAV at best endurance (E) & 16 \% of drag \\
Light fighter, sustained turn (G) & 13 \% of drag \\
Helicopter hover, 3 000 m ISA+20 (H) & 13.5 \% of rotor power \\
Helicopter hover, sea level & 17.6 \% of rotor power \\
Quadcopter hover (J) & 17 \% of rotor power \\
\bottomrule
\end{tabularx}
}
\end{table}

\hypertarget{rotors-and-propellers}{%
\subsection{Rotors and propellers}\label{rotors-and-propellers}}

For rotors and propellers the airfoil acts directly on the conversion of
shaft power into thrust, through the profile power term and through the
blade-element efficiency. These relations are developed in Sections V
and IX. The key point is that the airfoil drag here multiplies the
\emph{cube} of the local blade speed, so it is weighted strongly towards
the blade tip, where Mach number and compressibility effects are
largest.

\hypertarget{subsonic-jet-transports}{%
\section{Subsonic Jet Transports}\label{subsonic-jet-transports}}

\hypertarget{role-of-the-airfoil}{%
\subsection{Role of the airfoil}\label{role-of-the-airfoil}}

On a jet transport the wing is the largest single contributor to both
parasite and wave drag. Its sections are chosen to meet three coupled
requirements: low profile drag at cruise \(c_l\); a drag-divergence Mach
number comfortably above the cruise Mach number at that \(c_l\) (the
buffet boundary sets a further margin above that); and as much thickness
as possible, because thickness carries bending loads, houses fuel and
reduces structural weight. Raymer's statistical wing-weight relation for
transports, for example, makes wing weight vary roughly as
\((t/c)^{-0.4}\) {[}2{]}. The airfoil therefore touches the range
equation through \(L/D\) directly, through structural weight (and hence
\(W_i/W_f\)), and through the aspect ratio that a given thickness
permits.

\hypertarget{worked-example-a-profile-drag-and-fuel-burn}{%
\subsection{Worked example A: profile drag and fuel
burn}\label{worked-example-a-profile-drag-and-fuel-burn}}

Consider a generic single-aisle transport with wing reference area
\(S=122.6\) m², aspect ratio \(A=9.5\), Oswald factor \(e=0.82\) and
mid-cruise mass 66 t, cruising at \(M=0.78\) at 35 000 ft (10 668 m;
\(\rho=0.380\) kg/m³, \(V=231\) m/s, \(q=10\,154\) Pa). Its zero-lift
drag coefficient is taken as \(C_{D0}=0.0200\), of which the wing
profile drag is 0.0070 (35 \%), plus 10 counts of wave drag at cruise;
TSFC is 0.57 h\textsuperscript{−1}. Table IV gives the results.

At the mid-cruise \(C_L=0.520\), induced drag is \(C_{Di}=0.0110\) and
total drag \(C_D=0.0320\), giving \(L/D=16.2\). A 10 \% reduction in
wing profile drag (7 counts) raises \(L/D\) to 16.6. For a 2 500 km
cruise ending at 62 t, (9) gives a fuel burn of 6 897 kg for the
baseline and 6 739 kg for the improved wing: a saving of 159 kg or 2.3
\% per flight. Over 1 000 such flights---roughly one year of operation
for one aircraft---that is about 159 t of fuel and 500 t of
CO\textsubscript{2}, using the standard 3.16 kg CO\textsubscript{2} per
kg of jet fuel {[}23{]}.

The \emph{sensitivity} is the useful figure: each 1 \% of wing profile
drag is worth 0.22 \% of \(L/D\) and slightly more in fuel. Fig. 4 shows
that the gain applies across the cruise \(C_L\) range but is largest at
the low-\(C_L\) end of the polar, i.e.~early in a long flight at high
weight it is somewhat smaller, and late in the flight or at higher Mach
number it is larger.

\begin{figure}[!t]
\centering
\includegraphics[width=\columnwidth]{fig/f2_transport_LD.pdf}
\caption{Fig. 4. Lift-to-drag ratio of the generic transport, baseline and with
reduced wing profile drag. The dotted line marks the mid-cruise lift
coefficient.}
\end{figure}

\begin{table}[!t]
\centering
{\footnotesize TABLE IV\par\textsc{Worked example A: generic narrow-body transport at \(M\) 0.78,
FL350.}\par}\vspace{4pt}
{\footnotesize\setlength{\tabcolsep}{4pt}\renewcommand{\arraystretch}{1.12}
\begin{tabularx}{\columnwidth}{>{\raggedright\arraybackslash}Xr>{\raggedright\arraybackslash}X>{\raggedright\arraybackslash}X}
\toprule
\textbf{Quantity} & \textbf{Baseline} & \textbf{Wing \(c_d\) −10 \%} & \textbf{Laminar wing} \\
\midrule
Wing section \(c_d\) (counts) & 61 & 55 & 44 \\
\(C_{D0}\) & 0.0200 & 0.0193 & 0.0184 \\
\(C_D\) at \(C_L=0.520\) & 0.0320 & 0.0313 & 0.0304 \\
\(L/D\) & 16.22 & 16.59 & 17.08 \\
Fuel, 2 500 km (kg) & 6 897 & 6 739 & 6 534 \\
Saving per flight & --- & 159 kg (2.3 \%) & 363 kg (5.3 \%) \\
Saving per 1 000 flights & --- & 159 t fuel, 500 t CO\textsubscript{2} & 363 t fuel, 1 150 t CO\textsubscript{2} \\
\bottomrule
\end{tabularx}
}
\end{table}

\hypertarget{worked-example-b-laminar-flow}{%
\subsection{Worked example B: laminar
flow}\label{worked-example-b-laminar-flow}}

At the wing's mean-chord Reynolds number of \(2.6\times10^7\), the
section model (3) gives 61 counts for a fully turbulent 12 \% section
and 44 counts if transition is held at 40 \% chord on the upper surface
and 25 \% on the lower---a 29 \% reduction. Applying this to 80 \% of
the wing area (excluding the root, the fuselage-junction region and
areas behind slats and contaminated zones) raises \(L/D\) to 17.1 and
reduces trip fuel by 5.3 \% (Table IV).

That figure explains half a century of effort on transport laminar flow
{[}21{]}, {[}24{]}. The obstacles are equally clear. Wing sweep,
required for \(M_{DD}\), produces crossflow instability that limits
natural laminar flow to leading-edge sweeps below roughly 18--20° at
transport Reynolds numbers. Hybrid laminar flow control (HLFC) adds
leading-edge suction to damp crossflow and uses a favourable pressure
gradient further aft; it was flight tested on a Boeing 757 wing glove in
the early 1990s {[}21{]}. Airbus's BLADE demonstrator first flew
natural-laminar outer-wing panels on an A340 in September 2017, to study
the manufacturing tolerances needed in practice; the project's stated
aims---about half the wing friction drag and up to 5 \% lower
CO\textsubscript{2}---are consistent with the estimate above {[}25{]}.
In both cases the aerodynamic benefit has been established; the
remaining issues are surface quality, insect contamination at take-off
and landing, ice protection integrated with the suction system,
maintenance cost and the weight of the system {[}24{]}.

\hypertarget{transonic-drag-divergence}{%
\subsection{Transonic drag divergence}\label{transonic-drag-divergence}}

For a swept wing the drag-divergence Mach number can be estimated with
the Korn equation as extended to sweep {[}26{]}:

\[M_{DD} = \frac{\kappa_A}{\cos\Lambda} - \frac{t/c}{\cos^2\Lambda} - \frac{C_L}{10\cos^3\Lambda} \tag{12}\]

with \(\kappa_A\approx0.87\) for conventional (NACA 6-series type)
sections and up to about 0.95 for supercritical sections. The wave-drag
rise above the critical Mach number is commonly modelled by Lock's
fourth-power law, \(C_{Dw}=20(M-M_{cr})^4\) with
\(M_{cr}=M_{DD}-(0.1/80)^{1/3}\) {[}26{]}.

Fig. 5 shows (12) for \(C_L=0.55\). The supercritical airfoil's benefit
can be taken in three ways (worked example C): at a sweep of 25° and a
target \(M_{DD}=0.80\), the maximum thickness rises from 7.1 \% for a
conventional section to 14.3 \% for a supercritical one; alternatively,
at 11 \% thickness the sweep required falls from 35° to 13°; or, at
fixed sweep and thickness, \(M_{DD}\) rises from 0.75 to 0.84. Real
designs split the benefit, most often taking thickness (for lighter,
higher-aspect-ratio wings with more fuel volume) and a little Mach
number. This is the main way in which airfoil \emph{drag
characteristics}---specifically the wave-drag rise---affect transport
performance: indirectly, through the wing's geometry and weight, rather
than through a few counts of \(c_d\) at cruise {[}13{]}, {[}14{]},
{[}3{]}.

\begin{figure}[!t]
\centering
\includegraphics[width=\columnwidth]{fig/f3_korn.pdf}
\caption{Fig. 5. Drag-divergence Mach number from the Korn equation (12) at
\(C_L=0.55\) for conventional and supercritical sections at two sweep
angles. The dotted line marks \(M_{DD}=0.78\).}
\end{figure}

\hypertarget{mission-length-dependence}{%
\subsection{Mission-length dependence}\label{mission-length-dependence}}

Fig. 7 generalises example A. The fuel saving for a given profile-drag
reduction grows with mission length, because of the exponential in (9):
for a 10 \% reduction in wing profile drag the saving is 2.2 \% at 800
km and 2.5 \% at 6 000 km, and for a 50 \% reduction it is 11.1 \% and
12.2 \% respectively. In practice, on short sectors a larger share of
fuel is spent in climb, where higher \(C_L\) makes induced drag more
important, and the cruise-only figures overstate the benefit; on
long-haul sectors they understate it, because the lighter fuel load also
reduces induced drag.

\hypertarget{take-off-and-climb-the-high-lift-configuration}{%
\subsection{Take-off and climb: the high-lift
configuration}\label{take-off-and-climb-the-high-lift-configuration}}

Cruise is not the only phase in which section drag matters. At take-off,
a twin-engined transport must meet a minimum climb gradient of 2.4 \% in
the second segment with one engine inoperative. The gradient is
approximately \(\gamma = T/W - 1/(L/D)\), where \(L/D\) is that of the
take-off configuration---slats and flaps deployed, at a \(C_L\) well
above cruise---and is typically only 8--10. On a hot day or at a high
airport, this requirement, not runway length, often limits take-off
weight. The drag of the multi-element high-lift sections, including
slat-cove and flap-gap flows and the confluence of wakes, then enters
the payload directly.

For a generic case in which one engine delivers thrust equal to 13.5 \%
of a 78 t reference weight and the take-off \(L/D\) is 9.0, the gradient
requirement limits take-off mass to 77.9 t. Raising take-off \(L/D\) by
1 \% allows 640 kg more (0.82 \% per 1 \%); raising it by 5 \%---through
lower-drag flap settings, sealed slats or better high-lift
sections---allows 3.2 t (4.1 \%). The leverage is three to four times
larger than in cruise because the climb gradient is the small difference
between two larger numbers. This is consistent with the industry rule of
thumb, quoted by van Dam {[}27{]}, that a 1 \% improvement in take-off
\(L/D\) is worth several thousand pounds of payload on a large twin.
Fig. 6 shows the relation for three thrust levels.

\begin{figure}[!t]
\centering
\includegraphics[width=\columnwidth]{fig/f4b_takeoff.pdf}
\caption{Fig. 6. One-engine-inoperative climb-limited take-off mass versus
take-off \(L/D\) for a generic twin at three levels of available thrust
(2.4 \% second-segment gradient).}
\end{figure}

\begin{figure}[!t]
\centering
\includegraphics[width=\columnwidth]{fig/f4_fuel_saving.pdf}
\caption{Fig. 7. Cruise fuel saving versus reduction in wing profile drag for
three cruise ranges (generic transport of example A).}
\end{figure}

\hypertarget{turboprops-general-aviation-and-propellers}{%
\section{Turboprops, General Aviation and
Propellers}\label{turboprops-general-aviation-and-propellers}}

\hypertarget{airframe}{%
\subsection{Airframe}\label{airframe}}

Turboprop and general-aviation aircraft fly at Mach numbers of 0.2--0.6
with little or no sweep. Wave drag is negligible except near propeller
tips, so airfoil profile drag and the extent of laminar flow are the
dominant section characteristics. Low sweep and moderate Reynolds
numbers (\(3\times10^6\)--\(1.5\times10^7\)) make natural laminar flow
much more attainable than on jets, and NLF sections such as the
NLF(1)-0416 were developed specifically for this class {[}16{]}. Two
practical limits apply: the propeller slipstream, which washes a
significant fraction of the span with unsteady, turbulent flow and moves
transition forward periodically; and surface quality, particularly rivet
heads and joints in metallic structure.

For a generic regional turboprop with \(C_{D0}=0.026\) (wing share
0.009), \(A=12\), \(e=0.80\) and cruise \(C_L=0.50\), the cruise \(L/D\)
is 14.6; a 10 \% lower wing profile drag raises it to 15.0, a 2.7 \%
range improvement from (8) at constant fuel. The sensitivity is
\emph{larger} than for the jet (0.27 \% per \% versus 0.22 \%) because
the turboprop cruises further below its \(C_L^*\), where parasite drag
is a larger share of the total.

\hypertarget{worked-example-speed-of-a-light-aircraft}{%
\subsection{Worked example: speed of a light
aircraft}\label{worked-example-speed-of-a-light-aircraft}}

For a piston or turboprop aircraft at full power, the maximum level
speed follows from \(P\eta_p = DV\). At high speed, parasite drag
dominates and \(V_{\max}\propto(P\eta_p/C_{D0})^{1/3}\), so speed
responds only to the cube root of drag. A generic four-seat light
aircraft (1 150 kg, \(S=16.2\) m², \(A=7.4\), 150 kW, \(\eta_p=0.82\),
\(C_{D0}=0.030\) of which 0.012 is wing) has a maximum speed of 141 kt.
A 10 \% lower wing profile drag gives 143 kt; a natural-laminar wing
that removes 45 counts of wing drag gives 149 kt (+5.7 \%). Because
power required rises with the cube of speed, the same drag reduction is
worth more as \emph{fuel saved at a fixed cruise speed} than as
\emph{speed gained at fixed power}, which is how laminar-wing light
aircraft and business jets usually express their advantage. Fig. 8 makes
the point: at a fixed cruise speed of 130 kt, the laminar wing reduces
power and fuel flow by 13 \%, more than twice its percentage gain in top
speed.

\begin{figure}[!t]
\centering
\includegraphics[width=\columnwidth]{fig/f14_ga.pdf}
\caption{Fig. 8. Generic light aircraft: change in maximum level speed at full
power, and in power required at 130 kt, as the zero-lift drag
coefficient varies around the baseline of 0.030.}
\end{figure}

\hypertarget{propeller-blade-sections}{%
\subsection{Propeller blade sections}\label{propeller-blade-sections}}

The propeller is itself a set of rotating airfoils, and the term
\(\eta_p\) in (8) depends directly on blade section drag. For a blade
element with inflow angle \(\phi\) and drag-to-lift angle
\(\gamma=\tan^{-1}(c_d/c_l)\), the local efficiency, ignoring induced
losses, is {[}28{]}

\[\eta_{local} = \frac{\tan\phi}{\tan(\phi+\gamma)} \tag{13}\]

which is maximised at \(\phi = 45^\circ - \gamma/2\). Fig. 9 shows (13)
for section lift-to-drag ratios of 30--120. At the optimum, raising
\(c_l/c_d\) from 50 to 80 raises local efficiency from 0.961 to 0.975;
at \(\phi=20^\circ\)---typical of inboard sections in cruise and of the
whole blade at climb---the gain from \(c_l/c_d=30\) to 80 is from 0.905
to 0.962. Profile losses are thus a few per cent at the design point but
grow rapidly away from it, which is why propeller sections are chosen
for a wide low-drag range and why variable pitch, which keeps each
element near its best \(c_l/c_d\), was such an advance.

\begin{figure}[!t]
\centering
\includegraphics[width=\columnwidth]{fig/f8_prop_eta.pdf}
\caption{Fig. 9. Blade-element efficiency (13) versus inflow angle for several
section lift-to-drag ratios.}
\end{figure}

Near the tip the helical Mach number of a turboprop blade reaches
0.7--0.85 in cruise. Thin sections (4--6 \% at the tip) and, on advanced
propellers, sweep are used to delay drag divergence---the same transonic
physics as Section IV, appearing on a smaller scale. Because the
propeller's thrust and the airframe's drag combine multiplicatively in
(8), a 3 \% gain in \(\eta_p\) from better blade sections is worth as
much as a 3 \% gain in \(L/D\) from the wing; in the generic example the
two together would give about 6 \%.

\hypertarget{sailplanes-and-long-endurance-aircraft}{%
\section{Sailplanes and Long-Endurance
Aircraft}\label{sailplanes-and-long-endurance-aircraft}}

\hypertarget{sailplanes}{%
\subsection{Sailplanes}\label{sailplanes}}

The sailplane is the purest expression of airfoil-driven performance.
There is no engine; performance is measured by the glide ratio \(L/D\)
and the minimum sink rate, and wing profile drag is typically the
largest single drag term at the speeds used for inter-thermal cruise.
Sailplane airfoils by Wortmann, Eppler and Althaus, designed with codes
in the XFOIL tradition and verified in low-turbulence tunnels, hold
laminar flow over 60--95 \% of the chord on the lower surface and 50--70
\% on the upper, with turbulators to control the laminar separation
bubble near transition {[}4{]}.

Worked example D uses a generic 15 m span sailplane (\(S=10.5\) m²,
\(A=21.4\), \(e=0.92\), mass 450 kg) with a 15 \% thick laminar section
operating at a Reynolds number of about \(1.2\times10^6\). The section
model gives 60 counts with transition at 60 \%/70 \% chord. If insect
debris and rain trip transition to 8 \% on the upper surface and 40 \%
on the lower, the section drag rises to 95 counts. With 35 counts of
fuselage and tail drag added, the maximum glide ratio falls from 40.3 to
34.5 (−14 \%), the minimum sink rises from 0.65 to 0.70 m/s, and the
glide ratio at 100 km/h falls from 39.9 to 34.5 (Fig. 10). For a
competition pilot these differences are decisive, which is why
competition sailplanes carry bug-wipers and why the laminar bucket of a
sailplane section is made as wide as possible: it must cover the high
\(c_l\) of thermalling \emph{and} the low \(c_l\) of fast cruise, with
flaps used to shift the bucket between them.

\begin{figure}[!t]
\centering
\includegraphics[width=\columnwidth]{fig/f9_sailplane.pdf}
\caption{Fig. 10. Sink polar of the generic 15 m sailplane: clean laminar wing
and insect/rain-contaminated wing (worked example D).}
\end{figure}

\hypertarget{long-endurance-uavs}{%
\subsection{Long-endurance UAVs}\label{long-endurance-uavs}}

Medium- and high-altitude long-endurance (MALE/HALE) unmanned aircraft
share the sailplane's emphasis on high aspect ratio and laminar
sections, but they are optimised for endurance, \(C_L^{3/2}/C_D\),
rather than glide. As shown in Section III, endurance responds only to
the fourth root of \(C_{D0}\); the gain from lowering minimum section
drag is half as large, in percentage terms, as for range. The section
characteristics that matter more are the drag at \emph{high} \(c_l\)
(1.0--1.4 for loiter), the maximum lift coefficient and gentle stall,
and adequate thickness for a very slender, flexible, fuel-carrying wing.
At high altitude, low density cuts the chord Reynolds number to
\(5\times10^5\)--\(2\times10^6\), close to the range in which laminar
separation bubbles become significant, so HALE sections must combine
sailplane-like low drag with low-Reynolds-number robustness. Fig. 11
shows how quickly this happens. At constant equivalent airspeed, true
speed rises as \(1/\sqrt{\sigma}\) while density falls as \(\sigma\), so
the Reynolds number falls as \(\sqrt{\sigma}\) and also with the rising
viscosity-to-density ratio. A 1 m chord flying at 30 m/s equivalent
airspeed has \(Re=2.1\times10^6\) at sea level but only
\(8.1\times10^5\) at 18 km (\(\rho=0.121\) kg/m³); a 0.5 m chord at the
same condition enters the bubble-sensitive range.

For a generic HALE aircraft with \(A=25\), \(e=0.9\) and 80 counts of
non-wing drag, a 30 \% rise in wing section drag at altitude (from 80 to
104 counts, typical of a section designed at too high a Reynolds number)
reduces the endurance parameter \((C_L^{3/2}/C_D)_{\max}\) by 3.4 \%.
The figure is modest because of the fourth-root dependence in (10), but
it is in addition to the loss of maximum lift that usually accompanies
the same bubble, which forces the aircraft to loiter at a lower \(C_L\)
and costs more endurance still.

\begin{figure}[!t]
\centering
\includegraphics[width=\columnwidth]{fig/f15_hale_re.pdf}
\caption{Fig. 11. Chord Reynolds number versus altitude at constant equivalent
airspeed for three combinations of speed and chord (ISA). The shaded
band marks the range in which laminar separation bubbles strongly affect
section drag.}
\end{figure}

\hypertarget{small-aircraft-and-low-reynolds-number-flight}{%
\section{Small Aircraft and Low-Reynolds-Number
Flight}\label{small-aircraft-and-low-reynolds-number-flight}}

\hypertarget{the-laminar-separation-bubble}{%
\subsection{The laminar separation
bubble}\label{the-laminar-separation-bubble}}

Below a chord Reynolds number of about \(5\times10^5\), the favourable
behaviour of laminar boundary layers turns into a liability. A laminar
layer cannot withstand much adverse pressure gradient; it separates,
transition then occurs in the free shear layer, and the now-turbulent
flow may reattach as a \emph{laminar separation bubble}. The bubble
thickens the boundary layer downstream and adds pressure drag. As
Reynolds number falls toward \(5\times10^4\)--\(7\times10^4\) the bubble
lengthens and may fail to reattach at all, producing large drag
increases, loss of lift and hysteresis {[}8{]}, {[}10{]}. At these
Reynolds numbers section drag depends strongly on airfoil shape, surface
finish, free-stream turbulence and even propeller noise; thin, cambered
sections and trips or turbulators that force transition ahead of laminar
separation often outperform ``conventional'' thick sections {[}9{]}.

\hypertarget{worked-example-e-endurance-of-an-electric-mini-uav}{%
\subsection{Worked example E: endurance of an electric
mini-UAV}\label{worked-example-e-endurance-of-an-electric-mini-uav}}

Consider an electric fixed-wing UAV of 2.5 kg mass, \(S=0.5\) m²,
\(A=8\), \(e=0.8\) and chord 0.25 m, with 12 counts of non-wing parasite
drag, a 100 Wh battery of which 85 \% is usable, an overall propulsive
efficiency (motor, controller, propeller) of 0.55 and 3 W of avionics.
Its endurance is

\[E = \frac{\eta\,E_{batt}}{P_{req}/\eta_{prop} + P_{av}} \tag{14}\]

where \(P_{req}=\tfrac12\rho V^3 S C_D\). Two section polars are
compared at a Reynolds number of \(1.5\times10^5\): a well-chosen
low-Reynolds-number section with \(c_d = 0.010 + 0.012(c_l-0.6)^2\), and
a bubble-prone section with \(c_d = 0.014 + 0.030(c_l-0.4)^2\), which
represents a thicker airfoil designed for higher Reynolds numbers. At
the best-endurance condition (\(C_L\approx1.0\), \(V\approx9\) m/s), the
first section has 121 counts of drag and the aircraft needs 32.4 W,
giving 158 min endurance; the second has 226 counts, needs 37.4 W and
achieves 136 min---a 14 \% loss. Even though endurance responds only to
\(C_{D0}^{-1/4}\) for a \emph{flat} polar, the bubble penalty is
concentrated at high \(c_l\), exactly where an endurance aircraft flies,
so the loss is much larger than the minimum-drag figures alone suggest.
Fig. 12 shows the power-required curves: the two aircraft differ little
at high speed but by 15 \% at the slow speeds used for loiter.

\begin{figure}[!t]
\centering
\includegraphics[width=\columnwidth]{fig/f11_uav_power.pdf}
\caption{Fig. 12. Battery power required by the generic 2.5 kg UAV with a
well-chosen low-Reynolds-number section and with a bubble-prone section
(worked example E).}
\end{figure}

The lesson for small aircraft is that the airfoil should be chosen from
data at the right Reynolds number and the right \(c_l\); catalogue
polars at \(Re=10^6\) or above are misleading. XFOIL's \(e^N\)
transition model with a reduced \(N\) factor is generally reliable for
bubble prediction down to about \(Re=10^5\) when the flow reattaches
{[}1{]}, but wind-tunnel data such as those of Selig and co-workers
{[}9{]} remain the reference.

\hypertarget{combat-aircraft}{%
\section{Combat Aircraft}\label{combat-aircraft}}

Combat aircraft operate across the widest range of conditions---subsonic
cruise, transonic manoeuvre, supersonic dash---and their airfoils are
compromises among them. Three section characteristics dominate: wave
drag at supersonic speed, which demands very thin sections; the
lift-dependent drag at high \(c_l\) in transonic manoeuvre, which
depends on how much leading-edge suction the section can sustain; and
maximum lift and buffet onset.

\hypertarget{worked-example-f-supersonic-wave-drag}{%
\subsection{Worked example F: supersonic wave
drag}\label{worked-example-f-supersonic-wave-drag}}

In linearised supersonic flow the drag of a thin airfoil is {[}18{]}

\[c_d = \frac{4\alpha^2}{\beta} + \frac{2}{\beta}\left(\overline{\left(\frac{dy_u}{dx}\right)^2}+\overline{\left(\frac{dy_l}{dx}\right)^2}\right), \qquad \beta=\sqrt{M_\infty^2-1} \tag{15}\]

For a symmetric biconvex section of thickness ratio \(\tau\) the
thickness term is \(16\tau^2/(3\beta)\). At \(M=1.6\) (\(\beta=1.249\)),
this gives 38 counts for a 3 \% section, 68 counts for 4 \%, 107 for 5
\% and 154 for 6 \% (Fig. 13). For comparison, the turbulent skin
friction of both surfaces at \(Re=3\times10^7\) is about 51 counts.
Above about 4 \% thickness, \emph{wave drag of the section alone exceeds
its skin friction}, and it grows with the square of thickness. This is
why supersonic wings are thin---typically 3--5 \% for fighters---and why
the reduction in thickness has to be paid for in structural weight, fuel
volume and low-speed maximum lift. The lift-dependent term, which can be
written \(\beta C_L^2/4\) for a two-dimensional section, is also large;
wings are therefore swept so that the leading edge lies inside the Mach
cone (``subsonic leading edge''), allowing some leading-edge suction to
be recovered with rounded, cambered leading edges inboard.

\begin{figure}[!t]
\centering
\includegraphics[width=\columnwidth]{fig/f5_supersonic.pdf}
\caption{Fig. 13. Thickness wave drag of a biconvex section from linear theory
(15), compared with turbulent skin friction of both surfaces.}
\end{figure}

\hypertarget{worked-example-g-sustained-turn}{%
\subsection{Worked example G: sustained
turn}\label{worked-example-g-sustained-turn}}

Consider a generic light fighter of 10 t mass, \(S=40\) m² and available
thrust 60 kN at 5 000 m and \(M=0.8\) (\(q=24\,200\) Pa), with
\(C_{D0}=0.020\) (of which 0.008 is wing profile drag) and \(K=0.21\).
From (11) its sustained load factor is 4.41 at \(C_L=0.447\), a turn
rate of 9.42 °/s. A 10 \% reduction in wing profile drag raises this to
only 4.46 g and 9.51 °/s (+1 \%). A 10 \% reduction in \(K\)---which a
section with better leading-edge suction, a leading-edge flap or
variable camber can provide---raises it to 4.65 g and 9.96 °/s (+5.7
\%). For manoeuvre, therefore, the airfoil's \emph{drag rise with lift}
is roughly five times as important as its minimum drag. This is the
aerodynamic logic behind manoeuvre flaps, leading-edge root extensions
and the automatic camber scheduling of modern fighters. Fig. 14 extends
the comparison to reductions of up to 30 \%: the turn-rate gain from
reducing \(K\) grows steadily, whereas the gain from wing minimum drag
remains small throughout.

\begin{figure}[!t]
\centering
\includegraphics[width=\columnwidth]{fig/f12_fighter.pdf}
\caption{Fig. 14. Sustained turn rate of the generic light fighter at \(M\) 0.8
and 5 000 m as wing profile drag or the induced-drag factor \(K\) is
reduced (worked example G).}
\end{figure}

\hypertarget{transonic-buffet-and-high-incidence-drag}{%
\subsection{Transonic buffet and high-incidence
drag}\label{transonic-buffet-and-high-incidence-drag}}

Two further section characteristics shape combat aircraft performance.
First, in transonic manoeuvre the thin wing section reaches its
\emph{buffet onset} boundary---where shock-induced separation becomes
unsteady---at a lift coefficient well below its low-speed maximum.
Buffet limits usable \(C_L\) for tracking and weapon delivery, and the
drag rise that accompanies shock-induced separation cuts specific excess
power just as the pilot demands it. Sections and wing twist are
therefore tailored to delay buffet onset, and variable leading- and
trailing-edge flaps are scheduled with Mach number and incidence to keep
the section near its best drag-to-lift characteristic.

Second, at high incidence the drag of a thin section is dominated by the
loss of leading-edge suction: once the flow separates at the sharp
leading edge, the resultant force acts nearly normal to the chord and
the drag becomes roughly \(C_L\tan\alpha\). Combat aircraft with highly
swept leading edges turn this into an advantage through vortex lift, but
at the price of very high drag---useful for rapid deceleration but
costly for sustained turns. The airfoil's leading-edge radius and
camber, together with leading-edge devices, decide how far along this
path the aircraft travels before the drag penalty becomes prohibitive.

\hypertarget{helicopters}{%
\section{Helicopters}\label{helicopters}}

\hypertarget{hover}{%
\subsection{Hover}\label{hover}}

In hover the rotor's shaft power is the sum of induced and profile
power. With the empirical induced-power factor \(\kappa\) and a mean
blade drag coefficient \(\bar c_d\), momentum and blade-element theory
give {[}5{]}, {[}6{]}

\[P = \kappa\frac{T^{3/2}}{\sqrt{2\rho A}} + \frac{\sigma\bar c_d}{8}\rho A(\Omega R)^3 \tag{16}\]

where \(A=\pi R^2\) is disk area, \(\sigma\) the solidity and
\(\Omega R\) tip speed. The figure of merit, the ratio of ideal to
actual power, is

\[FM = \frac{C_T^{3/2}/\sqrt2}{\kappa C_T^{3/2}/\sqrt2 + \sigma\bar c_d/8} \tag{17}\]

The airfoil appears through \(\bar c_d\) evaluated at the mean blade
lift coefficient \(\bar c_l\approx6C_T/\sigma\). Here the classical
polar of Bailey for the NACA 0012,
\(c_d = 0.0087 - 0.0216\alpha + 0.4\alpha^2\) {[}29{]}, is used with an
empirical drag-rise term beyond the onset of stall. Fig. 15 shows \(FM\)
as a function of blade loading for three sections: an advanced
rotorcraft section (\(c_{d0}=0.0075\), later stall), the NACA 0012
baseline and an eroded or rough blade (\(c_{d0}=0.0115\), earlier
stall). At \(C_T/\sigma=0.08\) the figures of merit are 0.73, 0.71 and
0.67. The maximum values are 0.78, 0.75 and 0.73, reached at
\(C_T/\sigma\) of 0.14, 0.125 and 0.12 respectively. Better sections
thus give a higher peak \emph{and} move it to higher blade loading,
which allows lower solidity or lower tip speed for the same
thrust---both of which reduce profile power further and, through tip
speed, reduce noise.

\begin{figure}[!t]
\centering
\includegraphics[width=\columnwidth]{fig/f6_fm.pdf}
\caption{Fig. 15. Rotor figure of merit versus blade loading for three blade
sections (16)--(17). The fall at high \(C_T/\sigma\) represents onset of
blade stall.}
\end{figure}

\hypertarget{worked-example-h-hot-and-high-hover-payload}{%
\subsection{Worked example H: hot-and-high hover
payload}\label{worked-example-h-hot-and-high-hover-payload}}

Take a generic 5.5 t utility helicopter with a four-bladed main rotor of
6.6 m radius, \(\sigma=0.08\) (chord 0.415 m), \(\Omega R=215\) m/s and
\(\kappa=1.15\), hovering out of ground effect at 3 000 m, ISA+20 °C
(\(\rho=0.846\) kg/m³). Table V gives the results. The blade loading is
\(C_T/\sigma=0.126\), close to the figure-of-merit peak; induced power
is 947 kW and profile power 148 kW, 13.5 \% of the main-rotor total of 1
095 kW.

With engine power fixed at this value, the hover-capable gross mass
changes as follows: if the whole blade drag polar is 10 \% lower, it
rises by 47 kg; if it is 25 \% higher---typical of moderate leading-edge
erosion or sand and dust abrasion---it falls by 127 kg. Against a useful
load of perhaps 1.5--2 t, these are 2--3 \% and 6--8 \% respectively.
Small changes in the airfoil's drag therefore translate into substantial
changes in what a helicopter can lift at a hot, high landing site, and
erosion protection of the blade leading edge is a performance item as
much as a maintenance one.

Fig. 16 generalises the result for blade drag polars between 80 \% and
140 \% of the baseline. A perhaps surprising feature is that the effect
is \emph{larger at sea level} than at the hot-and-high condition: at sea
level the same helicopter hovers at a lower blade loading, profile power
is 17.6 \% of the total rather than 13.5 \%, and a 25 \% increase in
blade drag costs 182 kg rather than 127 kg. The airfoil's drag therefore
matters most where the rotor is lightly loaded, and its maximum lift and
stall drag rise matter most where it is heavily loaded.

\begin{figure}[!t]
\centering
\includegraphics[width=\columnwidth]{fig/f13_heli_payload.pdf}
\caption{Fig. 16. Change in hover-capable gross mass of the generic 5.5 t
helicopter at constant power as the blade drag polar is scaled, at sea
level and at 3 000 m ISA+20 °C.}
\end{figure}

\begin{table}[!t]
\centering
{\footnotesize TABLE V\par\textsc{Worked example H: generic 5.5 t helicopter in hot-and-high
hover.}\par}\vspace{4pt}
{\footnotesize\setlength{\tabcolsep}{4pt}\renewcommand{\arraystretch}{1.12}
\begin{tabularx}{\columnwidth}{>{\raggedright\arraybackslash}X>{\raggedright\arraybackslash}X}
\toprule
\textbf{Quantity} & \textbf{Value} \\
\midrule
Density, 3 000 m ISA+20 (kg/m³) & 0.846 \\
\(C_T\); \(C_T/\sigma\); mean \(\bar c_l\) & 0.0101; 0.126; 0.76 \\
Mean blade \(\bar c_d\) & 0.0129 \\
Induced power (kW) & 947 \\
Profile power (kW) & 148 (13.5 \%) \\
Main-rotor power (kW) & 1 095 \\
Figure of merit & 0.75 \\
Δ hover mass, blade polar −10 \% & +47 kg \\
Δ hover mass, blade polar +25 \% & −127 kg \\
\bottomrule
\end{tabularx}
}
\end{table}

\hypertarget{forward-flight}{%
\subsection{Forward flight}\label{forward-flight}}

In forward flight at advance ratio \(\mu=V/\Omega R\), main-rotor power
can be written

\[P = \kappa T v_i + \frac{\sigma\bar c_{d}}{8}\rho A(\Omega R)^3(1+4.65\mu^2) + \tfrac12\rho V^3 f \tag{18}\]

where \(v_i\) is the induced velocity and \(f\) the equivalent
flat-plate area of the fuselage and hub {[}5{]}. The induced term falls
with speed, the parasite term rises with its cube, and the profile term
rises with \(\mu^2\)---but (18) hides the most important airfoil effect.
On the advancing side the blade tip sees a Mach number

\[M_{AT} = \frac{\Omega R + V}{a} \tag{19}\]

which at 140 kt and sea level reaches 0.84 for the example rotor. A 12
\% thick section with a critical Mach number near 0.64 at zero lift is
then deep in its drag rise for a large part of each revolution.
Meanwhile, on the retreating side, the blade must generate its share of
lift at low dynamic pressure; local \(c_l\) climbs towards
\(c_{l,\max}\) and drag rises with lift until retreating-blade stall
sets a hard limit. The helicopter's maximum speed, its high-speed range
and its vibration level are thus set by the airfoil's drag behaviour at
\emph{both} ends of its operating envelope---high Mach number at low
\(c_l\), and low Mach number at high \(c_l\) {[}7{]}, {[}17{]}.

\hypertarget{worked-example-i-tip-section-and-high-speed-power}{%
\subsection{Worked example I: tip section and high-speed
power}\label{worked-example-i-tip-section-and-high-speed-power}}

To capture these effects the profile power was integrated over blade
radius and azimuth, with each element's Mach number from (19) and a
Lock-law compressibility drag increment. Two tip sections are compared
on the same sea-level rotor with \(f=1.8\) m²: a conventional 12 \%
thick section (\(\kappa_A=0.87\), \(c_{d0}=0.0087\),
\(M_{cr}\approx0.64\)) and a thin advanced tip section of 8 \% thickness
with supercritical-style design (\(\kappa_A=0.92\), \(c_{d0}=0.0080\),
\(M_{cr}\approx0.73\)).

Fig. 17 shows the result. In hover the profile power differs by only 8
\% (162 versus 149 kW). At 140 kt, however, it is 261 kW for the
conventional tip and 199 kW for the advanced tip---a 24 \% difference,
almost entirely from compressibility---and total main-rotor power falls
from 812 to 750 kW (−7.6 \%). Minimum power, at the best-endurance speed
of about 78 kt, differs by only 3 \%. The best specific range (distance
per unit energy) improves by 5.1 \%, and the best-range speed rises from
111 to 115 kt, because the advanced tip delays the compressibility
penalty that otherwise limits economical cruise. The airfoil's
\emph{drag-divergence behaviour} therefore buys high-speed capability,
while its \emph{minimum drag} buys hover and endurance.

\begin{figure}[!t]
\centering
\includegraphics[width=\columnwidth]{fig/f7_heli_fwd.pdf}
\caption{Fig. 17. Main-rotor power versus forward speed for two tip sections
(worked example I), with induced and parasite components. Dashed lines
show profile power.}
\end{figure}

This is the background to the families of dedicated rotorcraft airfoils
developed since the 1970s---the Boeing VR series, Sikorsky SC1095, ONERA
OA series and the RAE/Westland sections with the BERP tip, among others
{[}17{]}. These combine thin, high-\(M_{DD}\) outboard sections with
thicker, high-\(c_{l,\max}\) inboard sections, small pitching moments
(because blade torsional loads and control loads depend on \(c_m\)), and
tip planforms that further delay compressibility and stall. A Westland
Lynx with BERP blades set the absolute helicopter speed record of 400.87
km/h in 1986. Table VI summarises how the requirements change along the
blade.

\begin{table*}[!t]
\centering
{\footnotesize TABLE VI\par\textsc{Typical rotor-blade section requirements by radial station
(conventional helicopter at high speed).}\par}\vspace{4pt}
{\footnotesize\setlength{\tabcolsep}{4pt}\renewcommand{\arraystretch}{1.12}
\begin{tabularx}{\textwidth}{>{\raggedright\arraybackslash}X>{\raggedright\arraybackslash}Xl>{\raggedright\arraybackslash}X}
\toprule
\textbf{Blade region} & \textbf{Local \(M\) (hover / advancing)} & \textbf{\(c_l\) range} & \textbf{Section priorities} \\
\midrule
Root (0.2--0.5 \(R\)) & 0.15--0.35 / 0.3--0.6 & 0.4--1.2 & Thickness for structure, \(c_{l,\max}\), low drag at moderate \(Re\) \\
Mid-span (0.5--0.85 \(R\)) & 0.35--0.55 / 0.6--0.8 & 0.2--1.3 & High \(c_{l,\max}\) at \(M\approx0.4\), low \(c_m\), wide drag bucket \\
Tip (0.85--1.0 \(R\)) & 0.55--0.65 / 0.8--0.9 & −0.2--1.0 & Thin, high \(M_{DD}\) at low \(c_l\), low \(c_m\), compatible with tip
sweep \\
\bottomrule
\end{tabularx}
}
\end{table*}

\hypertarget{dynamic-stall-and-vibration}{%
\subsection{Dynamic stall and
vibration}\label{dynamic-stall-and-vibration}}

The retreating blade in high-speed or high-thrust flight undergoes
large, rapid incidence changes once per revolution. Under these
conditions the section does not stall quasi-statically: a leading-edge
vortex forms and is shed, producing transient overshoots in lift and
very large excursions in drag and nose-down pitching moment {[}5{]}. The
resulting torsional loads and vibrations, rather than the mean power,
usually limit thrust and speed. Static polars, even good ones, therefore
underestimate the drag that a rotor experiences near its envelope
boundary; sections with gentle trailing-edge stall and a high static
\(c_{l,\max}\) also tend to delay dynamic stall, which is a further
reason why the drag-rise behaviour of the section matters more for
rotorcraft than its minimum drag.

\hypertarget{tiltrotors-compounds-and-multicopters}{%
\section{Tiltrotors, Compounds and
Multicopters}\label{tiltrotors-compounds-and-multicopters}}

\hypertarget{tiltrotors}{%
\subsection{Tiltrotors}\label{tiltrotors}}

A tiltrotor's proprotor must be an efficient rotor in hover and an
efficient propeller in cruise. The two requirements pull blade twist and
section choice in opposite directions: hover favours moderate twist and
sections with low drag at \(c_l\approx0.5\)--0.7, while high-speed
cruise at large inflow angles favours much higher twist and thin,
high-\(M_{DD}\) tip sections. Equation (13) applies to cruise; Fig. 9
shows that at the large \(\phi\) of a tiltrotor in cruise the efficiency
penalty of a modest \(c_l/c_d\) is small \emph{if} each element is at
its best \(c_l\), so the difficulty is less the minimum drag than
keeping every section near its drag bucket in both modes {[}6{]},
{[}30{]}.

The tiltrotor's wing provides a second, less obvious airfoil-drag
effect. Because it must be stiff enough to avoid whirl flutter, the wing
is very thick (above 20 \% on current designs), which raises profile
drag and lowers \(M_{DD}\), contributing to the cruise-speed ceiling of
the configuration. In hover the wing sits in the rotor wake and
experiences a \emph{download} of order 10 \% of rotor
thrust---effectively the drag of the section at about 90° incidence.
Deflecting the flaperons to reduce the wing's projected area in the wake
reduces this download, and every percentage point of download recovered
adds roughly one percentage point of hover-capable gross mass {[}30{]}.
For a generic 20 t tiltrotor, reducing download from 10 \% to 7 \% of
rotor thrust raises hover-capable gross mass by about 3.3 \%, some 670
kg, which may be 10--15 \% of the useful payload. Here the relevant
``airfoil drag characteristic'' is the drag of the wing section at 90°
incidence, with and without flap deflection---far outside any
conventional polar, but decisive for this configuration.

\hypertarget{compound-helicopters}{%
\subsection{Compound helicopters}\label{compound-helicopters}}

Compound and lift-offset designs---with auxiliary propulsion and, in
some cases, wings---offload the rotor at high speed and slow it to keep
the advancing tip below drag divergence. Speeds of 250 kt and above have
been demonstrated, for example by the Eurocopter X3, which reached 255
kt in level flight in 2013. Slowing the rotor raises the advance ratio
to 0.5--0.8, so a large part of the retreating side lies in reversed
flow, where the section meets air from its trailing edge. Airfoil drag
characteristics in reverse flow, and at very high and very low local
Mach number on the same revolution, therefore become design drivers, and
section designs with blunter trailing edges or tailored camber are
studied specifically for these rotors.

\hypertarget{multicopters-and-evtol}{%
\subsection{Multicopters and eVTOL}\label{multicopters-and-evtol}}

Small multicopter drones combine low Reynolds number with high blade
loading. Blade chords of 15--25 mm at tip speeds of 60--100 m/s give
Reynolds numbers of \(3\times10^4\)--\(1.5\times10^5\), where profile
drag coefficients of 0.015--0.035 are typical {[}10{]}.

Worked example J considers a 2 kg quadcopter with four 254 mm (10 in)
two-bladed rotors (\(\sigma=0.10\), tip speed 80 m/s, \(\kappa=1.25\),
representative Reynolds number \(6.5\times10^4\)), a 74 Wh battery of
which 80 \% is usable, motor and controller efficiency 0.75 and 5 W of
avionics. Table VII gives the results. As \(\bar c_d\) rises from 0.015
to 0.035, the profile power of each rotor more than doubles (6.0 to 13.9
W), the figure of merit falls from 0.69 to 0.59 and hover endurance
falls from 14.7 to 12.5 min (−15 \%). For multicopters, hover is the
dominant flight condition and the airfoil's low-Reynolds-number drag is
therefore one of the few aerodynamic parameters with a first-order
effect on endurance, beside disk loading.

\begin{table}[!t]
\centering
{\footnotesize TABLE VII\par\textsc{Worked example J: 2 kg quadcopter hover.}\par}\vspace{4pt}
{\footnotesize\setlength{\tabcolsep}{4pt}\renewcommand{\arraystretch}{1.12}
\begin{tabular}{lrrrr}
\toprule
\textbf{Rotor \(\bar c_d\)} & \textbf{Profile power per rotor (W)} & \textbf{FM} & \textbf{Total electrical power (W)} & \textbf{Hover endurance (min)} \\
\midrule
0.015 & 6.0 & 0.69 & 242 & 14.7 \\
0.020 & 8.0 & 0.66 & 253 & 14.0 \\
0.025 & 9.9 & 0.64 & 263 & 13.5 \\
0.035 & 13.9 & 0.59 & 285 & 12.5 \\
\bottomrule
\end{tabular}
}
\end{table}

Larger electric VTOL aircraft for urban air mobility add a further
trade-off: noise regulations push tip speeds down, raising blade \(c_l\)
for the same thrust and moving the sections towards the high-\(c_l\) end
of their drag polar. Airfoils that combine low drag at high \(c_l\) with
benign stall are needed, much as for the endurance UAV of Section VII.

Fig. 18 illustrates the trade for a generic six-rotor, 2.2 t
lift-plus-cruise vehicle with rotors of 1.6 m radius and solidity 0.15
in hover. As tip speed falls from 230 m/s, profile power falls with its
cube and the figure of merit rises; but the blade loading rises at the
same time, and once the mean blade lift coefficient approaches stall the
drag rise reverses the trend. With a conventional section the best
figure of merit, 0.78, occurs at a tip speed of about 135 m/s and
performance collapses below about 130 m/s; with an advanced section
having lower drag and a later drag rise, the optimum moves to about 124
m/s with a figure of merit of 0.80. Since rotor noise falls steeply with
tip speed, the section's drag behaviour at high \(c_l\) translates
directly into how quiet the vehicle can be made without a performance
penalty---or, alternatively, how much solidity, and therefore blade
weight, must be added to reach a noise target.

\begin{figure}[!t]
\centering
\includegraphics[width=\columnwidth]{fig/f16_evtol.pdf}
\caption{Fig. 18. Hover figure of merit versus tip speed for a generic 2.2 t
six-rotor eVTOL lift rotor with two blade sections. The fall at low tip
speed is caused by blade stall as \(C_T/\sigma\) rises.}
\end{figure}

\begin{figure}[!t]
\centering
\includegraphics[width=\columnwidth]{fig/f10_quad.pdf}
\caption{Fig. 19. Hover endurance and figure of merit of the generic 2 kg
quadcopter versus rotor mean profile drag coefficient (worked example
J).}
\end{figure}

Fig. 19 shows the trend continuously: each 0.001 of \(\bar c_d\) costs
about 7 s of hover at this size, and the figure of merit falls almost
linearly. Small rotor blades are therefore worth making thin,
sharp-edged and smooth, and propeller catalogues that quote static
thrust at a single rotational speed hide most of this sensitivity.

\hypertarget{second-order-couplings-moment-trim-and-propulsion-integration}{%
\section{Second-Order Couplings: Moment, Trim and Propulsion
Integration}\label{second-order-couplings-moment-trim-and-propulsion-integration}}

\hypertarget{pitching-moment-and-trim-drag}{%
\subsection{Pitching moment and trim
drag}\label{pitching-moment-and-trim-drag}}

An airfoil is not described by its drag polar alone; its pitching moment
\(c_m\) also carries a drag cost. Supercritical and aft-loaded sections,
which obtain much of their lift near the trailing edge, have large
nose-down moments. On a conventional aircraft the horizontal tail must
balance this with a download, so the wing must carry more lift than the
aircraft weight. The result is \emph{trim drag}: extra induced drag on
the wing, plus the induced and profile drag of the loaded tail. Its size
depends on centre-of-gravity position and tail arm, and is commonly of
the order of one to a few per cent of cruise drag {[}3{]}. A section
chosen for 2 \% lower profile drag can therefore lose part of its
advantage at the aircraft level if its moment is much more negative.
Fuel-transfer systems that move the centre of gravity aft in cruise are,
in effect, a way to recover trim drag created partly by the choice of
airfoil.

The constraint is sharper elsewhere. Tailless aircraft and flying wings
need sections with near-zero or slightly positive moment (reflexed
camber), which generally costs minimum drag and maximum lift. Helicopter
blades, as noted in Section IX, need small \(c_m\) to limit torsional
and control loads; for decades this kept rotorcraft on symmetric
sections despite their higher drag, and advanced rotor airfoils use
trailing-edge tabs or reflex to keep moments small while recovering lift
and drag performance {[}17{]}.

\hypertarget{propulsion-integration}{%
\subsection{Propulsion integration}\label{propulsion-integration}}

Propellers and engines change the flow over the airfoil. A tractor
propeller slipstream raises the local dynamic pressure on the washed
part of the wing and makes the boundary layer periodically turbulent;
the profile drag of that part of the wing therefore rises, and
laminar-flow benefits are limited to the unwashed span. Distributed
electric propulsion, studied for example in NASA's X-57 programme,
places many small propellers along the leading edge to increase lift at
low speed, so that a smaller, more highly loaded wing---with lower
wetted area and profile drag in cruise---can be used. The airfoil must
then perform well both in a powerful slipstream at take-off and in clean
flow at cruise.

Engine placement also interacts with the section's pressure
distribution. Fujino's over-the-wing nacelle for the HondaJet was
positioned so that its interference with the wing's supersonic region
\emph{reduced} wave drag at cruise, and the same aircraft uses a
natural-laminar-flow wing section {[}31{]}. Boundary-layer ingestion, in
which an engine swallows the low-momentum fluid of a wing or fuselage
boundary layer, goes further: part of the energy lost to profile drag is
recovered as reduced propulsive power. In such configurations section
drag is no longer simply a loss to be minimised but a flow to be managed
jointly with the propulsor.

\hypertarget{case-studies}{%
\section{Case Studies}\label{case-studies}}

The worked examples above use generic parameters. The following short
case studies show the same principles in real designs; the descriptions
are qualitative and drawn from the cited open literature.

\hypertarget{laminar-flow-piston-fighters}{%
\subsection{Laminar-flow piston
fighters}\label{laminar-flow-piston-fighters}}

During the Second World War several designs adopted the new laminar-flow
sections in the hope of large drag reductions. In service, achieved
laminar extent was limited by manufacturing waviness, rivets, paint,
insect contamination and combat damage. Nevertheless, the sections'
higher critical Mach number and thinner, more favourable pressure
distributions proved valuable at the high diving speeds of the period
{[}11{]}, {[}12{]}---an early demonstration that the
\emph{compressibility} characteristics of a section can matter more in
practice than its minimum drag.

\hypertarget{natural-laminar-flow-on-a-business-jet}{%
\subsection{Natural laminar flow on a business
jet}\label{natural-laminar-flow-on-a-business-jet}}

The HondaJet combines a natural-laminar-flow wing section (SHM-1), a
natural-laminar-flow fuselage nose and over-the-wing engine mounts
{[}31{]}. To sustain laminar flow the wing uses an integrally machined
metal skin to minimise waviness, steps and gaps. Low sweep, moderate
Reynolds number and careful surface engineering---exactly the conditions
identified in Sections II and IV---allowed a laminar wing to reach
production on a jet aircraft, where transport-size wings with higher
sweep and Reynolds number have not yet done so.

\hypertarget{supersonic-transport}{%
\subsection{Supersonic transport}\label{supersonic-transport}}

Concorde's wing is a very thin, cambered and twisted ogival delta. Its
thinness limits supersonic wave drag, in line with (15); the price is a
low-speed regime in which lift comes largely from leading-edge vortices
at high incidence, with high induced drag, high take-off power and high
noise. The thin section, chosen for its supersonic drag characteristics,
effectively sets the aircraft's low-speed performance and
community-noise problem.

\hypertarget{high-speed-helicopter-rotor}{%
\subsection{High-speed helicopter
rotor}\label{high-speed-helicopter-rotor}}

The BERP rotor of the Westland Lynx combined advanced rotor sections
with a tip planform designed to delay both advancing-blade
compressibility and retreating-blade stall. It enabled the 1986 speed
record noted in Section IX and illustrates the conclusion of worked
example I: for a helicopter, the drag \emph{rise} of the tip sections
with Mach number and lift coefficient sets the speed envelope.

\hypertarget{rotorcraft-on-mars}{%
\subsection{Rotorcraft on Mars}\label{rotorcraft-on-mars}}

The Ingenuity Mars helicopter flew in an atmosphere with roughly one per
cent of Earth's sea-level density. Its rotor blades operated at chord
Reynolds numbers of only \(10^3\)--\(10^4\) while the low speed of sound
pushed tip Mach numbers towards 0.7 {[}32{]}, {[}33{]}. In this regime
viscous drag is a large fraction of rotor power, and studies showed that
thin cambered plates can outperform conventional airfoils, giving higher
maximum thrust and figure of merit {[}33{]}. This is the extreme limit
of the low-Reynolds-number trend shown in Sections VII and X.

\hypertarget{sailplane-evolution}{%
\subsection{Sailplane evolution}\label{sailplane-evolution}}

The move from wooden to glass-fibre sailplanes in the 1960s allowed
exact, smooth wing contours and made the full laminar potential of the
Wortmann and Eppler sections attainable. Maximum glide ratios rose from
about 30 to above 40 for standard-class gliders and above 60 for
open-class gliders, with laminar section drag, aspect ratio and
contamination control as the main levers {[}4{]}, {[}15{]}. Worked
example D reproduces the sensitivity behind this history.

\hypertarget{in-service-degradation-of-airfoil-drag}{%
\section{In-Service Degradation of Airfoil
Drag}\label{in-service-degradation-of-airfoil-drag}}

The airfoil drag that a vehicle experiences in service is rarely the
drag of the clean, as-designed section. Several mechanisms degrade it,
and their importance varies with vehicle type in a way that mirrors the
sensitivities found above.

\emph{Contamination and roughness.} Insect debris collected near the
ground, rain, dust and paint degradation all promote early transition.
For laminar sections the effect is large---worked example D showed a 14
\% loss of glide ratio---while for sections already turbulent the
penalty is the smaller increase of turbulent skin friction with
roughness height. The permissible roughness for laminar flow is a small
fraction of the local boundary-layer thickness, so surface standards for
laminar wings are far tighter than for conventional ones {[}21{]},
{[}24{]}.

\emph{Manufacturing tolerances.} Steps, gaps and waviness at panel
joints, fastener heads and leading-edge attachments can cause premature
transition. Experience from laminar-flow flight programmes shows that
achievable production tolerances, rather than aerodynamic design, now
limit how much laminar flow can be relied upon {[}21{]}, {[}25{]}.

\emph{Icing.} Ice accretion changes the leading-edge geometry. Even thin
roughness-type ice raises drag considerably and lowers maximum lift;
horn and ridge ice can raise section drag several-fold and cause early
stall {[}34{]}. For fixed-wing aircraft this reduces climb and stall
margins; for helicopters it raises rotor torque and can exhaust the
power margin in hover. Ice protection systems are certified against
these performance losses.

\emph{Erosion.} Helicopter rotor blades, propellers and fan blades
operating in sand, dust and rain suffer leading-edge erosion that
roughens and reshapes the section. Worked example H showed that a 25 \%
increase in blade drag costs about 127 kg of hot-and-high hover
capability on a 5.5 t helicopter. Protective erosion strips (nickel,
titanium) and polyurethane tapes are therefore performance items as well
as maintenance items, and their own effect on section shape near the
leading edge must be considered.

\emph{Aeroelastic deformation.} Flexible wings and blades twist and bend
under load, changing the local \(c_l\) at which each section operates
and moving it within or out of its drag bucket. This is significant for
very high-aspect-ratio wings and for rotor blades, whose elastic twist
varies around the azimuth.

\hypertarget{predicting-airfoil-drag}{%
\section{Predicting Airfoil Drag}\label{predicting-airfoil-drag}}

Reliable performance prediction depends on reliable section drag, and
the methods available differ in cost and fidelity.

\emph{Handbook and empirical methods} such as (2)--(3), Hoerner's data
{[}20{]} and component build-up methods {[}2{]}, {[}22{]} are adequate
for conceptual design and for trend studies such as those in this paper.
They do not capture pressure-gradient effects on transition, laminar
separation bubbles or detailed compressibility behaviour.

\emph{Viscous--inviscid interaction codes} couple a panel or Euler
solution of the outer flow to an integral boundary-layer formulation
with \(e^N\) transition prediction. XFOIL {[}1{]} for incompressible and
weakly compressible flow, and ISES/MSES {[}35{]} for transonic flow, are
the standard tools; they predict minimum drag and drag-bucket behaviour
of attached-flow sections typically within a few drag counts, and
laminar separation bubbles reasonably well down to \(Re\approx10^5\).
Their limitations---at and beyond maximum lift, in massive separation,
at very low Reynolds number and in unsteady conditions such as the
oscillating incidence of a rotor blade---are the regions where
performance is often most sensitive to the airfoil (Sections VII--IX).
Extending such codes, for example with roughness and contamination
models, improved post-stall correlations or rotor-relevant unsteady
corrections, is a productive direction for improvement.

\emph{RANS CFD with transition models} (for example the
correlation-based \(\gamma\)--\(Re_\theta\) model {[}36{]}) handles
complex geometry, shocks and separation, but predicted drag remains
sensitive to grid and turbulence modelling, and transition modelling in
particular is still under development. \emph{Large-eddy and direct
simulation} resolve bubbles and transition physics but remain research
tools for sections.

\emph{Wind-tunnel and flight measurement}, by wake-rake momentum deficit
in low-turbulence tunnels and by flight-test wake surveys, remain the
reference for section drag, especially at low Reynolds number {[}9{]}
and for laminar sections whose behaviour depends on free-stream
turbulence. A sound design process uses fast codes to explore,
higher-fidelity CFD to check, and test data to calibrate. Table VIII
summarises the trade-offs.

\begin{table*}[!t]
\centering
{\footnotesize TABLE VIII\par\textsc{Methods for obtaining section drag data.}\par}\vspace{4pt}
{\footnotesize\setlength{\tabcolsep}{4pt}\renewcommand{\arraystretch}{1.12}
\begin{tabularx}{\textwidth}{>{\raggedright\arraybackslash}Xl>{\raggedright\arraybackslash}X>{\raggedright\arraybackslash}X}
\toprule
\textbf{Method} & \textbf{Cost per polar} & \textbf{Strengths} & \textbf{Weak regimes} \\
\midrule
Handbook build-up, eqs. (2)--(3) & Negligible & Trends, conceptual sizing & Pressure-gradient effects, bubbles, stall, transonic \\
Viscous--inviscid (XFOIL, MSES) & Seconds & Minimum drag, drag bucket, transition, bubbles to \(Re\approx10^5\),
optimisation & Post-stall, massive separation, very low \(Re\), unsteady \\
RANS with transition model & Hours & Shocks, separation, 3-D geometry, high-lift systems & Transition and separation sensitivity to models and grids \\
LES / DNS & Days to weeks & Bubble and transition physics & Cost; design use limited \\
Wind tunnel (wake rake) & Days, high cost & Reference data, real surfaces, low-\(Re\) & Tunnel turbulence, wall corrections, \(Re\) scaling \\
Flight test & Very high & True operating environment, contamination & Limited conditions, instrumentation \\
\bottomrule
\end{tabularx}
}
\end{table*}

\hypertarget{airfoil-selection-and-multi-point-design}{%
\section{Airfoil Selection and Multi-Point
Design}\label{airfoil-selection-and-multi-point-design}}

The results above suggest that the classical practice of selecting an
airfoil for a single design point---minimum drag at the cruise \(c_l\),
say---is inadequate for most vehicles. A better formulation weights the
section drag by the time, energy or power that each operating condition
accounts for. For a fixed-wing aircraft a natural objective is

\[J = \sum_i w_i\, c_d(c_{l,i}, Re_i, M_i) \tag{20}\]

where the conditions \(i\) cover climb, cruise at several weights,
loiter and descent, and the weights \(w_i\) reflect fuel burned (or
energy used) in each, multiplied by the local sensitivity of fuel to
drag derived in Section III. For a rotor blade the natural weighting is
by the cube of the local velocity, because profile power at a blade
element is proportional to \(U^3 c_d\); the objective is then an
integral over radius and azimuth,

\[J_{rotor} = \int_0^{2\pi}\!\!\int_{r_0}^{R} U(r,\psi)^3\, c_d\big(c_l(r,\psi), Re, M(r,\psi)\big)\, c(r)\,dr\,d\psi \tag{21}\]

which automatically emphasises the advancing tip and penalises
compressibility drag there, as worked example I showed.

Such objectives must be accompanied by constraints that encode the other
roles of the section: minimum thickness (for structure, fuel or
stiffness), maximum lift and stall behaviour (for take-off, landing,
retreating blade or manoeuvre), pitching moment (for trim drag or
control loads), buffet margin, and robustness to roughness. The last is
particularly important: a section optimised for the clean condition can
be fragile, losing much of its advantage when transition is forced. A
robust formulation evaluates the drag at both free and fixed transition
and penalises the difference, or optimises the expected drag over a
distribution of surface conditions.

Viscous--inviscid codes such as XFOIL are well suited to these
multi-point formulations because each evaluation takes a fraction of a
second, allowing thousands of candidate shapes to be assessed with
gradient-based or evolutionary optimisers. The weak points, as noted in
Section XIV, are high-lift, post-stall and strongly unsteady conditions,
so multi-point optimisations usually verify the final shapes with RANS
or wind-tunnel data in those regions. Extending open codes so that the
mission-weighted objectives (20)--(21) and the robustness constraints
can be expressed directly would make them considerably more useful for
vehicle-level design.

\hypertarget{synthesis-and-design-guidance}{%
\section{Synthesis and Design
Guidance}\label{synthesis-and-design-guidance}}

Table IX collects the results of the worked examples. The final column
normalises each result, where meaningful, to the performance change per
1 \% change of section drag.

\begin{table*}[!t]
\centering
{\footnotesize TABLE IX\par\textsc{Summary of worked examples.}\par}\vspace{4pt}
{\footnotesize\setlength{\tabcolsep}{4pt}\renewcommand{\arraystretch}{1.12}
\begin{tabularx}{\textwidth}{l>{\raggedright\arraybackslash}X>{\raggedright\arraybackslash}X>{\raggedright\arraybackslash}X>{\raggedright\arraybackslash}X}
\toprule
\textbf{Ex.} & \textbf{Vehicle} & \textbf{Change in airfoil characteristic} & \textbf{Performance effect} & \textbf{Per 1 \% section drag} \\
\midrule
A & Jet transport & Wing profile drag −10 \% & Trip fuel −2.3 \% (2 500 km) & ≈0.23 \% fuel \\
B & Jet transport & Laminar wing: section \(c_d\) −29 \% on 80 \% of area & Trip fuel −5.3 \% & ≈0.23 \% fuel \\
C & Jet transport & Supercritical vs conventional section & \(t/c\) 7.1 \% → 14.3 \% at \(M_{DD}\) 0.80, or sweep 35° → 13° & Indirect: weight, \(A\) \\
--- & Jet transport, take-off & Take-off \(L/D\) +1 \% / +5 \% (high-lift sections) & OEI-limited take-off mass +640 kg / +3.2 t & ≈0.8 \% mass per 1 \% \(L/D\) \\
--- & Light aircraft & Wing profile drag −10 \%; laminar wing (−45 counts) & \(V_{\max}\) +1.4 \%; +5.7 \% & ≈0.14 \% speed \\
--- & Turboprop & Wing profile drag −10 \% & Range +2.7 \% & ≈0.27 \% range \\
--- & Propeller & Blade \(c_l/c_d\) 50 → 80 & Local efficiency +1.5 points at optimum & --- \\
D & Sailplane & Contamination: section \(c_d\) +58 \% & \((L/D)_{\max}\) 40.3 → 34.5 (−14 \%) & ≈0.25 \% glide ratio \\
E & Mini-UAV & Bubble-prone section (\(c_d\) +87 \% at loiter \(c_l\)) & Endurance −14 \% & ≈0.16 \% endurance \\
F & Supersonic & Section \(t/c\) 4 \% → 3 \% at \(M\) 1.6 & Wave drag 68 → 38 counts & \(\propto (t/c)^2\) \\
G & Fighter & Wing profile drag −10 \%; or \(K\) −10 \% & Turn rate +1.0 \%; or +5.7 \% & ≈0.1 \%; \(K\) ≈0.6 \% \\
H & Helicopter, hover & Blade polar −10 \% / +25 \% & +47 kg / −127 kg hover mass & ≈0.09 \% of gross mass \\
I & Helicopter, 140 kt & Thin high-\(M_{DD}\) tip section & Rotor power −7.6 \%; best specific range +5.1 \% & Mach-driven \\
J & Quadcopter & Rotor \(\bar c_d\) 0.015 → 0.035 & FM 0.69 → 0.59; endurance −15 \% & ≈0.11 \% endurance \\
\bottomrule
\end{tabularx}
}
\end{table*}

A complementary view is the \emph{value of one drag count}, which
designers use to set priorities among competing refinements. Table X
gives it for the generic vehicles.

\begin{table}[!t]
\centering
{\footnotesize TABLE X\par\textsc{Value of one drag count (\(\Delta C_D\) or \(\Delta \bar c_d\)
= 0.0001) for the generic vehicles.}\par}\vspace{4pt}
{\footnotesize\setlength{\tabcolsep}{4pt}\renewcommand{\arraystretch}{1.12}
\begin{tabularx}{\columnwidth}{>{\raggedright\arraybackslash}X>{\raggedright\arraybackslash}X>{\raggedright\arraybackslash}X}
\toprule
\textbf{Vehicle (example)} & \textbf{One count of} & \textbf{Value} \\
\midrule
Jet transport (A) & Aircraft \(C_D\) & 22.7 kg fuel per 2 500 km flight (0.33 \%) \\
Regional turboprop & Aircraft \(C_D\) & 0.29 \% range \\
Sailplane (D) & Aircraft \(C_D\) & +0.21 in \((L/D)_{\max}\) \\
Light fighter (G) & Aircraft \(C_D\) & +0.005 g sustained load factor \\
Helicopter, hot-and-high (H) & Blade mean \(\bar c_d\) & 3.7 kg hover-capable mass \\
Quadcopter (J) & Rotor mean \(\bar c_d\) & 0.7 s hover endurance \\
\bottomrule
\end{tabularx}
}
\end{table}

Fig. 20 presents the same normalised sensitivities graphically. Only one
bar---the fighter's response to the lift-dependent drag factor
\(K\)---stands well clear of the others, which reinforces the point that
for manoeuvre-limited aircraft the \emph{shape} of the section polar at
high lift matters more than its minimum.

\begin{figure}[!t]
\centering
\includegraphics[width=\columnwidth]{fig/f17_sens.pdf}
\caption{Fig. 20. Normalised sensitivity of the main performance metric of each
generic vehicle to the drag quantity governed by the airfoil (per cent
change in metric per per cent change in drag).}
\end{figure}

\hypertarget{economic-interpretation}{%
\subsection{Economic interpretation}\label{economic-interpretation}}

Sensitivities become decisions when they are converted into value. For
the generic transport of example A, one drag count costs 22.7 kg of fuel
per 2 500 km flight, or about 23 t of fuel and 72 t of
CO\textsubscript{2} per 1 000 flights. At an illustrative fuel price of
USD 800 per tonne, that is of the order of USD 18 000 per aircraft per
year for each count, and nearly USD 2 million per year for a fleet of
100 aircraft. It is this arithmetic that justifies the engineering
effort spent on a few counts of wing drag, and also sets the ceiling on
what a laminar-flow system may cost in weight, maintenance and
complexity. For the helicopter of example H, the corresponding currency
is payload at the hottest and highest landing site in the mission: 3.7
kg per count of blade drag, or 127 kg for the erosion case, may decide
whether a mission can be flown with one aircraft or needs two. For small
drones the currency is flight time, and a minute of hover endurance can
be the difference between completing a survey line and returning early.

\hypertarget{robustness-of-the-results}{%
\subsection{Robustness of the results}\label{robustness-of-the-results}}

Because the worked examples use generic parameters, it is worth checking
how sensitive the conclusions are to those parameters. Table XI repeats
worked example A---the fuel saving from a 10 \% reduction in wing
profile drag---with the main assumptions varied one at a time. The
result ranges from 2.0 \% to 2.6 \%. The largest single influence is the
wing's share of zero-lift drag, which the airfoil designer does not
control; aspect ratio, cruise lift coefficient and TSFC have smaller
effects. The \emph{magnitude} of the sensitivity, about a quarter of a
per cent of fuel per per cent of wing profile drag, is therefore robust;
its precise value for a particular aircraft needs that aircraft's drag
breakdown. Similar checks on the other examples give the same picture:
the numbers move by 10--25 \%, but the ranking of what matters for each
vehicle class does not change.

\begin{table}[!t]
\centering
{\footnotesize TABLE XI\par\textsc{Robustness of worked example A: cruise fuel saving for a 10 \%
reduction in wing profile drag (2 500 km).}\par}\vspace{4pt}
{\footnotesize\setlength{\tabcolsep}{4pt}\renewcommand{\arraystretch}{1.12}
\begin{tabularx}{\columnwidth}{>{\raggedright\arraybackslash}Xlr}
\toprule
\textbf{Assumption varied} & \textbf{Value} & \textbf{Fuel saving (\%)} \\
\midrule
Baseline & --- & 2.30 \\
Aspect ratio & 8.5 / 10.5 & 2.21 / 2.37 \\
Wing share of \(C_{D0}\) & 30 \% / 40 \% & 1.97 / 2.63 \\
TSFC (h\textsuperscript{−1}) & 0.53 / 0.62 & 2.29 / 2.31 \\
Cruise \(C_L\) & −10 \% / +10 \% & 2.46 / 2.14 \\
\bottomrule
\end{tabularx}
}
\end{table}

Five general conclusions follow.

\begin{enumerate}
\def\labelenumi{\arabic{enumi}.}
\item
  \textbf{The limiting performance regime determines which airfoil
  characteristic matters.} Range-limited aircraft cruising below
  \(C_L^*\) (jets, turboprops) respond to minimum profile drag and its
  extent; endurance-limited aircraft (MALE UAVs, small drones) respond
  to drag at high \(c_l\); manoeuvre-limited aircraft respond to the
  lift-dependent drag rise; Mach-limited aircraft (transports,
  helicopter advancing blades, supersonic aircraft) respond to drag
  divergence and wave drag.
\item
  \textbf{Sensitivities are of similar magnitude---about 0.1--0.3 \% of
  the performance metric per 1 \% of section drag---across very
  different vehicles}, because the wing or rotor profile drag is
  typically 15--40 \% of the relevant total. The exceptions are regimes
  where the airfoil's drag \emph{rise} is triggered---transonic,
  supersonic and stall-limited---where effects are far larger and more
  abrupt.
\item
  \textbf{Indirect effects often exceed direct ones.} For the transport,
  the supercritical airfoil's value lies more in the thicker, lighter,
  higher-aspect-ratio wing it permits than in the drag counts it saves
  at cruise. For the helicopter, a better section allows lower solidity
  or tip speed, saving more profile power and reducing noise.
\item
  \textbf{Laminar flow is the largest single opportunity and the most
  fragile.} It can halve section skin friction, but its value depends on
  maintaining it in service against contamination, manufacturing
  tolerances and sweep.
\item
  \textbf{Scale matters.} As vehicles shrink---small UAVs, multicopters,
  rotor blade tips at low Reynolds number---the airfoil moves into a
  regime where drag depends strongly on transition behaviour and laminar
  separation, and high-Reynolds-number data become misleading.
\end{enumerate}

For designers, these conclusions suggest a practical procedure. Identify
which regime limits the vehicle's key mission metric. Determine the
section \(c_l\), \(Re\) and \(M\) range actually encountered, including
azimuthal variation for rotors and climb and loiter conditions for fixed
wings. Select or design sections for that range, not for a single design
point. Then estimate the sensitivity of the mission metric to section
drag with relations of the kind used here, so that aerodynamic,
structural and operational trade-offs---laminar flow versus maintenance,
thickness versus weight, tip speed versus noise---can be made on a
common basis.

\hypertarget{conclusions}{%
\section{Conclusions}\label{conclusions}}

This paper has traced how the drag characteristics of the
two-dimensional airfoil section propagate into the performance of
fixed-wing aircraft and rotorcraft. Using a common set of simple,
calibrated relations---a mixed-boundary-layer section-drag model, the
parabolic polar, the Breguet relations, the Korn and Lock transonic
relations, linear supersonic theory and blade-element momentum
theory---ten worked examples were evaluated for transports, turboprops
and propellers, sailplanes, small UAVs, combat aircraft, helicopters and
multicopters.

The main findings are that a 10 \% change in airfoil profile drag
typically changes range, endurance or hover capability by 1--3 \%; that
a laminar transport wing could save about 5 \% of trip fuel; that
in-service contamination and erosion can cost more performance than many
design refinements gain; and that, for rotorcraft and high-speed
aircraft, the drag \emph{rise} with Mach number and lift coefficient
matters more than minimum drag, because it sets the performance
boundary. The analysis also shows where better prediction tools are most
valuable: near maximum lift, at low Reynolds number, in transonic flow
and in the unsteady environment of a rotor. These are the natural
targets for improving open airfoil analysis codes.

Several research directions follow from this analysis. \emph{Robust
laminar flow}---sections and surfaces that retain much of their laminar
benefit when contaminated or manufactured to realistic
tolerances---would move the transport result of worked example B closer
to practice. \emph{Passive and active skin-friction reduction}, such as
riblets for turbulent regions and suction or plasma actuation for
transition control, attacks the same term in (2). \emph{Variable
geometry}, from conventional flaps to continuous trailing-edge morphing,
allows a section to keep its drag bucket aligned with a moving operating
point, which Sections VI, IX and X showed to be as valuable as lowering
the bucket itself. \emph{Rotor-specific design}, including sections for
reverse flow on slowed compound rotors and for very low Reynolds numbers
on small and planetary rotors, remains comparatively immature. Finally,
\emph{mission-weighted, robust multi-point design tools} of the kind
outlined in Section XV would make it easier to turn these airfoil-level
gains into the vehicle-level benefits that this paper has quantified.

\hypertarget{acknowledgement}{%
\section*{Acknowledgement}\label{acknowledgement}}
\addcontentsline{toc}{section}{Acknowledgement}

The author acknowledges the use of Claude, a large language model
developed by Anthropic, for assistance in drafting the text, organising
the literature and preparing the worked-example calculations and
figures. The author reviewed and takes responsibility for the content.
The worked examples use generic, representative parameters and are not
models of any specific aircraft.

\hypertarget{appendix-a-parameters-of-the-worked-examples}{%
\section*{Appendix A: Parameters of the Worked
Examples}\label{appendix-a-parameters-of-the-worked-examples}}
\addcontentsline{toc}{section}{Appendix A: Parameters of the Worked
Examples}

\begin{table*}[!t]
\centering
{\footnotesize TABLE XII\par\textsc{Parameters used in the worked examples (all generic and
representative).}\par}\vspace{4pt}
{\footnotesize\setlength{\tabcolsep}{4pt}\renewcommand{\arraystretch}{1.12}
\begin{tabularx}{\textwidth}{l>{\raggedright\arraybackslash}X}
\toprule
\textbf{Ex.} & \textbf{Main parameters} \\
\midrule
A, B & \(S\)=122.6 m², \(A\)=9.5, \(e\)=0.82, mid-cruise mass 66 t, \(M\)=0.78,
FL350, \(C_{D0}\)=0.0200 (wing 0.0070), \(C_{Dw}\)=0.0010, TSFC 0.57
h\textsuperscript{−1}, 2 500 km cruise ending at 62 t; laminar:
transition 40 \%/25 \%, \(Re\)=2.6×10⁷, 80 \% of wing treated \\
C & Korn equation, \(C_L\)=0.55, \(\Lambda\)=25°, \(\kappa_A\)=0.87/0.95 \\
Turboprop & \(C_{D0}\)=0.026 (wing 0.009), \(A\)=12, \(e\)=0.80, \(C_L\)=0.50 \\
D & \(b\)=15 m, \(S\)=10.5 m², \(e\)=0.92, 450 kg, \(t/c\)=15 \%,
\(Re\)=1.2×10⁶, clean transition 60 \%/70 \%, contaminated 8 \%/40 \%,
other drag 0.0035 \\
E & 2.5 kg, \(S\)=0.5 m², \(A\)=8, \(e\)=0.8, other drag 0.012, 100 Wh (85
\% usable), \(\eta\)=0.55, avionics 3 W \\
F & Biconvex section, linear theory, \(M\)=1.6 \\
G & 10 t, \(S\)=40 m², \(T\)=60 kN, 5 000 m, \(M\)=0.8, \(C_{D0}\)=0.020
(wing 0.008), \(K\)=0.21 \\
H & 5.5 t, \(R\)=6.6 m, 4 blades, \(\sigma\)=0.08, \(\Omega R\)=215 m/s,
\(\kappa\)=1.15, 3 000 m ISA+20, Bailey NACA 0012 polar with stall drag
rise \\
I & As H at sea level, \(f\)=1.8 m², blade-element profile power over radius
and azimuth, Lock compressibility increment \\
J & 2 kg, 4 rotors \(R\)=0.127 m, \(\sigma\)=0.10, \(\Omega R\)=80 m/s,
\(\kappa\)=1.25, 74 Wh (80 \% usable), \(\eta\)=0.75, avionics 5 W \\
\bottomrule
\end{tabularx}
}
\end{table*}

\hypertarget{appendix-b-analytical-sensitivities}{%
\section*{Appendix B: Analytical
Sensitivities}\label{appendix-b-analytical-sensitivities}}
\addcontentsline{toc}{section}{Appendix B: Analytical Sensitivities}

The numerical results can be checked against simple logarithmic
derivatives of the governing relations. Writing
\(\epsilon_x^y = \partial\ln y/\partial\ln x\) for the sensitivity of a
metric \(y\) to a drag quantity \(x\):

\emph{Lift-to-drag ratio at fixed \(C_L\).} Since \(L/D=C_L/C_D\), a
change in the wing profile drag contribution \(C_{D,w}\) gives
\(\epsilon^{L/D}_{C_{D,w}} = -C_{D,w}/C_D\). For example A this is
\(-0.0070/0.0320 = -0.22\), matching the numerical result.

\emph{Fuel at fixed range.} From (9), with \(x = Rc/(V\,L/D)\),
\(\epsilon^{W_{fuel}}_{L/D} = -x e^x/(e^x-1)\). For example A,
\(x=0.105\) and the factor is \(-1.05\); combined with the previous
result, each 1 \% of wing profile drag costs 0.23 \% of trip fuel.

\emph{Optimum-point metrics.} From (6) and (10),
\(\epsilon^{(L/D)_{\max}}_{C_{D0}} = -\tfrac12\) and
\(\epsilon^{(C_L^{3/2}/C_D)_{\max}}_{C_{D0}} = -\tfrac14\); from
\(P\eta_p = DV\) at high speed,
\(\epsilon^{V_{\max}}_{C_{D0}} \approx -\tfrac13\).

\emph{Sustained turn.} From (11),
\(\partial n_{sus}/\partial C_{D0} = -n_{sus}/[2(T/qS - C_{D0})]\) and
\(\epsilon^{n_{sus}}_{K} = -\tfrac12\). For example G, one drag count
changes \(n_{sus}\) by 0.005, and 10 \% of \(K\) changes it by about 5
\%, as computed.

\emph{Climb-limited take-off weight.} With \(\gamma = T/W - (L/D)^{-1}\)
held at the regulatory minimum,
\(\epsilon^{W}_{L/D} = (L/D)^{-1}/[\gamma + (L/D)^{-1}]\), which is 0.82
for \(\gamma=0.024\) and \(L/D=9\).

\emph{Hover at fixed power.} From (16) with \(P = aT^{3/2} + b\bar c_d\)
held constant, \(\epsilon^{T}_{\bar c_d} = -\tfrac23 P_0/P_i\). For
example H, \(P_0/P_i = 148/947\), so each 1 \% of blade drag costs about
0.10 \% of hover thrust; the full calculation gives a slightly smaller
figure because \(\bar c_d\) itself falls as \(C_T\) falls.

\emph{Multicopter endurance.} With electrical power proportional to
\(P_i + P_0\) plus a fixed avionics load,
\(\epsilon^{E}_{\bar c_d} \approx -P_0/(P_i+P_0+P_{av}')\), where
\(P_{av}'\) is the avionics load referred to the shaft. For example J at
\(\bar c_d=0.02\), this is about \(-0.17\) near the baseline, consistent
with the 0.7 s per count of Table X.

These expressions show compactly why the sensitivities in Table IX
cluster between 0.1 and 0.3: the relevant drag share (\(C_{D,w}/C_D\),
\(P_0/P\), and so on) is typically 15--35 \%, and the exponent
multiplying it is of order one.

\hypertarget{appendix-c-step-by-step-calculation-for-worked-example-a}{%
\section*{Appendix C: Step-by-Step Calculation for Worked Example
A}\label{appendix-c-step-by-step-calculation-for-worked-example-a}}
\addcontentsline{toc}{section}{Appendix C: Step-by-Step Calculation for
Worked Example A}

The following steps reproduce the baseline and improved results of Table
IV.

\begin{enumerate}
\def\labelenumi{\arabic{enumi}.}
\tightlist
\item
  \emph{Atmosphere.} At 10 668 m in the ISA, \(T=218.8\) K,
  \(\rho=0.3796\) kg/m³ and \(a=296.5\) m/s. At \(M=0.78\), \(V=231.3\)
  m/s and \(q=\tfrac12\rho V^2=10\,154\) Pa.
\item
  \emph{Lift coefficient.}
  \(C_L = W/(qS) = 66\,000\times9.807/(10\,154\times122.6) = 0.520\).
\item
  \emph{Induced drag.} \(K = 1/(\pi\times9.5\times0.82)=0.0409\), so
  \(C_{Di}=KC_L^2=0.0110\).
\item
  \emph{Total drag and \(L/D\).}
  \(C_D = 0.0200 + 0.0010 + 0.0110 = 0.0320\) and
  \(L/D = 0.520/0.0320 = 16.22\).
\item
  \emph{Fuel.} With \(c = 0.57/3600 = 1.583\times10^{-4}\) s⁻¹,
  \(x = Rc/(V\,L/D) = 2.5\times10^6\times1.583\times10^{-4}/(231.3\times16.22) = 0.1055\),
  and from (9) \(W_{fuel}/g = 62\,000\,(e^{x}-1) = 6\,897\) kg.
\item
  \emph{Improved wing.} Reducing wing profile drag by 10 \% lowers
  \(C_{D0}\) by 0.0007 to 0.0193, giving \(C_D=0.0313\), \(L/D=16.59\),
  \(x=0.1032\) and 6 739 kg of fuel: a saving of 159 kg (2.3 \%).
\item
  \emph{Laminar wing.} The section model gives 61.1 counts turbulent and
  43.6 counts with transition at 40 \%/25 \% at \(Re=2.6\times10^7\), a
  28.6 \% reduction; applied to 80 \% of the wing it lowers \(C_{D0}\)
  by 0.0016 to 0.0184, giving \(L/D=17.08\) and 6 534 kg of fuel.
\end{enumerate}

\hypertarget{appendix-d-reproducibility}{%
\section*{Appendix D:
Reproducibility}\label{appendix-d-reproducibility}}
\addcontentsline{toc}{section}{Appendix D: Reproducibility}

All numerical results and figures in this paper were produced by a
single Python script (calc\_figs.py, using NumPy and Matplotlib) that
implements equations (1)--(21) and the blade-element integration of
worked example I, writes every computed value to a results file, and
draws each figure. The paper itself is assembled from Markdown sources
by a build script that numbers equations, figures, tables and references
and produces both the web version and this two-column PDF. Readers
wishing to apply the method to a specific aircraft need only replace the
generic parameters of Table XII with that aircraft's data and drag
breakdown.

\hypertarget{references}{%
\section*{References}\label{references}}
\addcontentsline{toc}{section}{References}

\begingroup\footnotesize\begin{list}{}{\setlength{\leftmargin}{2.2em}\setlength{\labelwidth}{2em}\setlength{\labelsep}{0.3em}\setlength{\itemsep}{1pt}\setlength{\parsep}{0pt}}

\item[{[1]}] M. Drela, ``XFOIL: An analysis and design system for low
Reynolds number airfoils,'' in \emph{Low Reynolds Number Aerodynamics},
T. J. Mueller, Ed., Lecture Notes in Engineering, vol.~54. Berlin,
Germany: Springer, 1989, pp.~1--12.

\item[{[2]}] D. P. Raymer, \emph{Aircraft Design: A Conceptual Approach}, 6th
ed.~Reston, VA, USA: AIAA, 2018.

\item[{[3]}] E. Torenbeek, \emph{Advanced Aircraft Design}. Chichester, U.K.:
Wiley, 2013.

\item[{[4]}] F. Thomas, \emph{Fundamentals of Sailplane Design}. College
Park, MD, USA: College Park Press, 1999.

\item[{[5]}] J. G. Leishman, \emph{Principles of Helicopter Aerodynamics},
2nd ed.~Cambridge, U.K.: Cambridge Univ. Press, 2006.

\item[{[6]}] W. Johnson, \emph{Rotorcraft Aeromechanics}. Cambridge, U.K.:
Cambridge Univ. Press, 2013.

\item[{[7]}] R. W. Prouty, \emph{Helicopter Performance, Stability, and
Control}. Malabar, FL, USA: Krieger, 1995.

\item[{[8]}] P. B. S. Lissaman, ``Low-Reynolds-number airfoils,'' \emph{Annu.
Rev.~Fluid Mech.}, vol.~15, pp.~223--239, 1983.

\item[{[9]}] M. S. Selig, J. F. Donovan, and D. B. Fraser, \emph{Airfoils at
Low Speeds} (Soartech 8). Virginia Beach, VA, USA: H. A. Stokely, 1989.

\item[{[10]}] T. J. Mueller and J. D. DeLaurier, ``Aerodynamics of small
vehicles,'' \emph{Annu. Rev.~Fluid Mech.}, vol.~35, pp.~89--111, 2003.

\item[{[11]}] I. H. Abbott and A. E. von Doenhoff, \emph{Theory of Wing
Sections}. New York, NY, USA: Dover, 1959.

\item[{[12]}] L. K. Loftin, Jr., \emph{Quest for Performance: The Evolution
of Modern Aircraft}, NASA SP-468, 1985.

\item[{[13]}] R. T. Whitcomb and L. R. Clark, ``An airfoil shape for
efficient flight at supercritical Mach numbers,'' NASA TM X-1109, 1965.

\item[{[14]}] C. D. Harris, ``NASA supercritical airfoils: A matrix of
family-related airfoils,'' NASA TP-2969, 1990.

\item[{[15]}] R. Eppler, \emph{Airfoil Design and Data}. Berlin, Germany:
Springer, 1990.

\item[{[16]}] D. M. Somers, ``Design and experimental results for a
natural-laminar-flow airfoil for general aviation applications,'' NASA
TP-1861, 1981.

\item[{[17]}] W. G. Bousman, ``Airfoil design and rotorcraft performance,''
in \emph{Proc. American Helicopter Society 58th Annual Forum}, Montréal,
QC, Canada, Jun.~2002, pp.~208--232.

\item[{[18]}] J. D. Anderson, Jr., \emph{Fundamentals of Aerodynamics}, 6th
ed.~New York, NY, USA: McGraw-Hill, 2017.

\item[{[19]}] H. Schlichting and K. Gersten, \emph{Boundary-Layer Theory},
9th ed.~Berlin, Germany: Springer, 2017.

\item[{[20]}] S. F. Hoerner, \emph{Fluid-Dynamic Drag}. Brick Town, NJ, USA:
Hoerner Fluid Dynamics, 1965.

\item[{[21]}] R. D. Joslin, ``Aircraft laminar flow control,'' \emph{Annu.
Rev.~Fluid Mech.}, vol.~30, pp.~1--29, 1998.

\item[{[22]}] O. Gur, W. H. Mason, and J. A. Schetz, ``Full-configuration
drag estimation,'' \emph{J. Aircraft}, vol.~47, no. 4, pp.~1356--1367,
2010.

\item[{[23]}] International Civil Aviation Organization, \emph{ICAO Carbon
Emissions Calculator Methodology}. Montréal, QC, Canada: ICAO.

\item[{[24]}] D. M. Bushnell, ``Aircraft drag reduction---A review,''
\emph{Proc. Inst. Mech. Eng. G, J. Aerosp. Eng.}, vol.~217, no. 1,
pp.~1--18, 2003.

\item[{[25]}] Airbus, ``Airbus' BLADE laminar flow wing demonstrator makes
first flight,'' press release, Toulouse, France, 26 Sep.~2017.

\item[{[26]}] W. H. Mason, \emph{Configuration Aerodynamics}, course notes,
ch.~7 (Transonic aerodynamics of airfoils and wings). Blacksburg, VA,
USA: Virginia Polytechnic Inst. and State Univ.

\item[{[27]}] C. P. van Dam, ``The aerodynamic design of multi-element
high-lift systems for transport airplanes,'' \emph{Prog. Aerosp. Sci.},
vol.~38, no. 2, pp.~101--144, 2002.

\item[{[28]}] B. W. McCormick, \emph{Aerodynamics, Aeronautics, and Flight
Mechanics}, 2nd ed.~New York, NY, USA: Wiley, 1995.

\item[{[29]}] F. J. Bailey, Jr., ``A simplified theoretical method of
determining the characteristics of a lifting rotor in forward flight,''
NACA Rep.~716, 1941.

\item[{[30]}] M. D. Maisel, D. J. Giulianetti, and D. C. Dugan, \emph{The
History of the XV-15 Tilt Rotor Research Aircraft: From Concept to
Flight}, NASA SP-2000-4517, 2000.

\item[{[31]}] M. Fujino, ``Design and development of the HondaJet,'' \emph{J.
Aircraft}, vol.~42, no. 3, pp.~755--764, 2005.

\item[{[32]}] J. Balaram \emph{et al.}, ``Mars Helicopter technology
demonstrator,'' in \emph{Proc. AIAA Atmospheric Flight Mechanics Conf.}
(SciTech Forum), Kissimmee, FL, USA, Jan.~2018, AIAA 2018-0023.

\item[{[33]}] W. J. F. Koning, E. A. Romander, and W. Johnson, ``Low Reynolds
number airfoil evaluation for the Mars Helicopter rotor,'' in
\emph{Proc. AHS International 74th Annual Forum}, Phoenix, AZ, USA, May
2018.

\item[{[34]}] M. B. Bragg, A. P. Broeren, and L. A. Blumenthal,
``Iced-airfoil aerodynamics,'' \emph{Prog. Aerosp. Sci.}, vol.~41, no.
5, pp.~323--362, 2005.

\item[{[35]}] M. Drela and M. B. Giles, ``Viscous-inviscid analysis of
transonic and low Reynolds number airfoils,'' \emph{AIAA J.}, vol.~25,
no. 10, pp.~1347--1355, 1987.

\item[{[36]}] R. B. Langtry and F. R. Menter, ``Correlation-based transition
modeling for unstructured parallelized computational fluid dynamics
codes,'' \emph{AIAA J.}, vol.~47, no. 12, pp.~2894--2906, 2009.

\end{list}\endgroup

\end{document}
